\documentclass[11pt,a4paper]{article}
\usepackage{amsmath,amssymb}
\usepackage{booktabs}
\usepackage{array}
\usepackage{longtable}
\usepackage{geometry}
\usepackage[hidelinks]{hyperref}
\geometry{margin=2.5cm}

\title{\textbf{GRB Afterglow Detection Rates: Implementation Reference}\\
\large Model equations as implemented in \texttt{grb\_detections\_claude}}
\author{}
\date{}

\begin{document}
\maketitle
\thispagestyle{empty}

% -------------------------------------------------------
\section*{Summary}

The code computes the expected annual detection rate $R_{\rm det}$ [yr$^{-1}$] of gamma-ray burst (GRB) optical afterglows as a function of survey strategy parameters $(N_{\rm exp},\,t_{\rm cad})$.
The calculation combines a Blandford--McKee ISM afterglow model with a Euclidean population of GRBs and a background-limited survey instrument.

The central insight is that detectability is governed by \emph{two competing horizons}.
The \textbf{Euclidean horizon} $q_{\rm Euc}$ is the maximum off-axis viewing angle at which a burst at the calibration distance $D_{\rm euc}$ is brighter than the survey flux limit.
The \textbf{cadence horizon} $q_i$ is the maximum off-axis angle whose afterglow peak arrives early enough to be caught within $i$ cadence intervals.
Comparing these two angles against three characteristic afterglow angles ($q_{\rm dec}$, $q_j$, $q_{\rm nr}$) divides the strategy space into seven piecewise regimes A1--A7, each with its own analytic rate formula.
The surface $\log_{10}R_{\rm det}(N_{\rm exp},t_{\rm cad})$ is computed on a dense grid and displayed interactively.

% -------------------------------------------------------
\section{Parameters}

{\scriptsize
\renewcommand{\arraystretch}{1.2}
\setlength{\tabcolsep}{3pt}
\begin{longtable}{
  >{\raggedright\arraybackslash}p{1.8cm}
  >{\raggedright\arraybackslash}p{2.9cm}
  >{\raggedright\arraybackslash}p{2.9cm}
  >{\raggedright\arraybackslash}p{1.6cm}
  >{\raggedright\arraybackslash}p{5.6cm}}
\toprule
Symbol & Default & First value & Unit & Meaning \\
\midrule
\endhead
\multicolumn{5}{l}{\textit{Physical model} (\texttt{AfterglowPhysicalParams})} \\
$p$ & 2.2 & 2.5* & -- & Electron energy distribution index \\
$E_{k,{\rm iso}}$ & $10^{53}$ & (unchanged) & erg & Isotropic kinetic energy \\
$n_0$ & 1.0 & (unchanged) & cm$^{-3}$ & ISM ambient number density \\
$\theta_j$ & 0.1 & (unchanged) & rad & Jet half-opening angle \\
$\Gamma_0$ & $10^{2.5}$ & (unchanged) & -- & Initial bulk Lorentz factor \\
$\nu$ & $5\times10^{14}$ & (unchanged) & Hz & Observing frequency (optical V-band) \\
$z_{\rm Euc}$ & 2.0 & (unchanged) & -- & Alt.\ cosmology redshift (reserved, unused by default) \\
$H_0$ & 70 & (unchanged) & km/s/Mpc & Alt.\ cosmology Hubble constant (reserved, unused) \\
$\Omega_m$ & 0.3 & (unchanged) & -- & Alt.\ cosmology matter density (reserved, unused) \\
\midrule
\multicolumn{5}{l}{\textit{Microphysics} (\texttt{MicrophysicsParams})} \\
$\varepsilon_e$ & 0.1 & (unchanged) & -- & Electron energy fraction \\
$\varepsilon_B$ & $10^{-4}$ & 0.01* & -- & Magnetic field energy fraction \\
$z$ & 0.0 & (unchanged) & -- & Redshift factor input (reserved, unused by default) \\
\texttt{include\_\allowbreak redshift\_\allowbreak factors} & False & (unchanged) & -- & Toggle for redshift corrections (reserved, unused) \\
\midrule
\multicolumn{5}{l}{\textit{Cosmological model}} \\
$\mathcal{R}$ & 260 & (unchanged) & Gpc$^{-3}$\,yr$^{-1}$ & Volumetric GRB rate density \\
$D_{\rm euc}$ & 1.404$\times10^{28}$\,cm $\approx$ 4.55\,Gpc & 1.63$\times10^{28}$\,cm $\approx$ 5.28\,Gpc* & cm & Euclidean calibration distance \\
$R_{\rm int}$ & $\frac{4\pi}{3}\mathcal{R}D_{\rm euc}^3\approx1.03\times10^5$ & $\approx1.60\times10^5$* & yr$^{-1}$ & Total GRBs/yr in calibration volume \\
\midrule
\multicolumn{5}{l}{\textit{Survey instrument} (\texttt{SurveyTelescopeParams}, \texttt{SurveyDesignParams})} \\
$F_{\rm ref}$ & $10^{-4.68}\approx2.09\times10^{-5}$ & (unchanged) & Jy & Limiting flux anchor at $t_{\rm exp,ref}$ \\
$\alpha$ & 0.5 & (unchanged) & -- & Flux--exposure exponent (background-limited) \\
$\Omega_{\rm exp}$ & 47 & (unchanged) & deg$^2$ & Field of view per exposure \\
$t_{\rm exp,ref}$ & 30 & (unchanged) & s & Reference exposure time for the $F_{\rm ref}$ anchor \\
$\Omega_{\rm srv,max}$ & 27500\textsuperscript{\dag} & (unchanged) & deg$^2$ & Maximum surveyable sky area \\
$t_{\rm night}$ & 36000\,s\,($=10$\,hr) & (unchanged) & s & Astronomical night duration \\
\midrule
\multicolumn{5}{l}{\textit{Survey strategy} (free variables)} \\
$N_{\rm exp}$ & -- & -- & -- & Exposures per cadence cycle \\
$t_{\rm cad}$ & -- & -- & s & Cadence interval (optical mode: integer days only, Sec.~\ref{sec:optical_constraints}) \\
$f_{\rm live}$ & 0.1\textsuperscript{\dag} & (unchanged) & -- & Wall-clock live fraction — the non-optical mode's budget parameter; in optical mode it is derived, $f_{\rm live}=f_{\rm eff}\,f_{\rm night}$ (Sec.~\ref{sec:schedule}) \\
$f_{\rm eff}$ & 0.48\textsuperscript{\dag} & (introduced 2026-08) & -- & Usable fraction of the night window — the optical mode's budget parameter (Sec.~\ref{sec:schedule}); undefined without a night, so each mode exposes its own slider \\
$N_v$ & 1 & (introduced 2026-08) & -- & Visits per visit-night (night schedule, Sec.~\ref{sec:schedule}); 1 = legacy model \\
$\Delta t_v$ & 2\,hr (UI) & (introduced 2026-08) & s & Intra-night visit spacing (night schedule); inert at $N_v=1$ \\
$t_{\rm OH}$ & 0 & (unchanged) & s & Fixed overhead per exposure \\
$i$ & 2\textsuperscript{\ddag} & 10*\textsuperscript{\ddag} & -- & Required detections per GRB (a pipeline property, decoupled from $N_v$) \\
\midrule
\multicolumn{5}{l}{\textit{Constraints} (detection-rate filters)} \\
$q_{\rm min}$ & 0 & 0 (introduced 2026-05-15) & -- & Minimum viewing-angle parameter; 0 = no filter \\
$D_{\rm min}$ & 0 & 0 (introduced 2026-05-15) & Gpc & Minimum detection distance; 0 = no filter \\
$s_{\rm fade}$ & 0 & 0 (introduced 2026-07, as \texttt{s\_min}) & mag/day & Minimum fading-rate identification cut; 0 = no filter \\
$s_{\rm rise}$ & 0 & 0 (introduced 2026-07) & mag/day & Minimum rise-rate identification cut; 0 = no filter \\
\bottomrule
\end{longtable}
}

\noindent\textbf{* Changed since introduction} (\texttt{grb\_detect/params.py}'s history, first
committed as \texttt{684f9c5}, 2026-02-20): $p$ 2.5$\to$2.2 (this update); $\varepsilon_B$
$0.01\to10^{-3.4}$ (\texttt{76eaa31}, 2026-07-20) $\to10^{-4}$ (a later update); $D_{\rm euc}$
5.28$\to$4.55\,Gpc, carrying $R_{\rm int}$ with it (\texttt{76eaa31}, 2026-07-20); $i$ 10$\to$2
(\texttt{76eaa31}, 2026-07-20). All other rows are unchanged from the first commit.

\noindent\textsuperscript{\dag} The interactive UI intentionally ships different initial slider
values than these library defaults: $\Omega_{\rm srv,max}=41253$\,deg$^2$ (full sky) and
$f_{\rm live}=0.2$. The library defaults above are what a bare constructor call
(\texttt{SurveyDesignParams()}, \texttt{SurveyTelescopeParams()}) still uses.

\noindent\textsuperscript{\ddag} $i$ has no library-level default anywhere in
\texttt{grb\_detect/params.py} (always a required argument); the values shown are the
interactive UI's initial slider position instead --- first the original Dash app's
\texttt{dcc.Slider} default, now \texttt{build\_standalone.py}'s slider default.

\noindent\textbf{Note on $R_{\rm int}$:} The \texttt{AfterglowPhysicalParams} dataclass bakes $R_{\rm int}$ in at class-load time from the default $\mathcal{R}$ and $D_{\rm euc}$ values.
When constructing the class with non-default $\mathcal{R}$ or $D_{\rm euc}$, always recompute $R_{\rm int}$ explicitly and pass it to the constructor:
\[
R_{\rm int} = \tfrac{4\pi}{3}\,\mathcal{R}\,D_{\rm euc,\,Gpc}^3
\]

% -------------------------------------------------------
\section{Afterglow Model (ISM)}

\subsection{Deceleration and jet-break times}

\noindent Blandford--McKee ISM dynamics (thesis Eq.~30):
\begin{equation}
t_{\rm dec} = 27\,\zeta\,E_{52}^{1/3}\,n_0^{-1/3}\,\Gamma_{2.5}^{-8/3}\;[\text{s}],
\qquad \zeta = \tfrac{1}{3}\;\text{(Euclidean approximation, }z\to0\text{)}
\end{equation}
where $E_{52}=E_{k,{\rm iso}}/10^{52}$\,erg and $\Gamma_{2.5}=\Gamma_0/10^{2.5}$.

\begin{equation}
t_j = t_{\rm dec}\,(\Gamma_0\,\theta_j)^{8/3}
\end{equation}

\subsection{Viewing-angle parameterisation}

\noindent Define normalized viewing angle $q = \theta_{\rm obs}/\theta_j$ and $\tilde{q}=q-1$.
Three characteristic boundaries separate the afterglow into phases:
\begin{equation}
q_{\rm dec} = 1 + (\Gamma_0\,\theta_j)^{-1}, \qquad q_j = 2, \qquad q_{\rm nr} = \frac{\sqrt{2}}{\theta_j}
\end{equation}
Phase~I ($t < t_{\rm dec}$) is excluded by design. Phase~II: $q_{\rm dec}\le q<q_j$. Phase~III: $q\ge q_j$.

\noindent\textbf{Beaming identity.} The normalization $q_{\rm nr}=\sqrt{2}/\theta_j$ is chosen so that
\[
\frac{\theta_j^2}{2}\,q_{\rm nr}^2 = 1,
\]
i.e.\ the integral over all observable solid angles $\int_0^{q_{\rm nr}}\theta_j^2 q\,dq = 1$ exactly.
This ensures that in regime A1 (saturation), $R_{\rm A1} = f_\Omega\,R_{\rm int}$ without any additional beaming prefactor.
For all other regimes the factor $\frac{1}{2}\theta_j^2 q_{\rm max}^2 < 1$ represents the fraction of GRBs detectable at viewing angles $q\le q_{\rm max}$.

\subsection{Off-axis peak flux (PLS~G, $\nu_m<\nu<\nu_c$)}

\noindent On-axis reference flux at $t_{\rm dec}$ (thesis Eq.~24):
\begin{equation}
F_{\rm dec}^{\rm G} = \frac{0.461}{1000}\,(p-0.04)\,e^{2.53p}\;
\bar{\varepsilon}_e^{\,p-1}\,\varepsilon_B^{(p+1)/4}\,
n_0^{1/2}\,E_{52}^{(p+3)/4}\,
t_{\rm dec,d}^{-3(p-1)/4}\,D_{28}^{-2}\,\nu_{14}^{-(p-1)/2}
\;[\text{Jy}]
\end{equation}
where $\bar{\varepsilon}_e=\varepsilon_e(p-2)/(p-1)$, $t_{\rm dec,d}=t_{\rm dec}/{\rm day}$, $D_{28}=D_{\rm euc}/10^{28}$\,cm, $\nu_{14}=\nu/10^{14}$\,Hz. Valid for $p>2$.

\noindent Off-axis flux exponents (PLS~G):
\[
a_{\rm II}^{\rm G} = p-1, \qquad a_{\rm III}^{\rm G} = p.
\]
These follow from the afterglow temporal decay index $F\propto t^{-3(p-1)/4}$ in Phase~II (combined with $t_p\propto\tilde{q}^{8/3}$), and $F\propto t^{-p}$ after the jet break (Phase~III, with $t_p\propto\tilde{q}^{2}$).

\subsection{Off-axis reference flux (PLS~H, $\nu>\nu_c$)}

\noindent For PLS~H (thesis Eq.~25), the on-axis reference flux is:
\begin{equation}
F_{\rm dec}^{\rm H} = \frac{0.855}{1000}\,(p-0.98)\,e^{1.95p}\;
\bar{\varepsilon}_e^{\,p-1}\,\varepsilon_B^{(p-2)/4}\,
E_{52}^{(p+2)/4}\,
t_{\rm dec,d}^{-(3p-2)/4}\,D_{28}^{-2}\,\nu_{14}^{-p/2}
\;[\text{Jy}]
\end{equation}
Key differences from PLS~G: no $n_0$ dependence; $\varepsilon_B$ exponent $(p-2)/4$; $E_{52}$ exponent $(p+2)/4$; time exponent $-(3p-2)/4$; $\nu$ exponent $-p/2$.

\noindent Off-axis flux exponents (PLS~H):
\[
a_{\rm II}^{\rm H} = \frac{3p-2}{3}, \qquad a_{\rm III}^{\rm H} = p.
\]
The Phase~II exponent follows from temporal slope $-(3p-2)/4$ combined with $t_p\propto\tilde{q}^{8/3}$ in Phase~II.

\subsection{Off-axis peak flux scaling}

\noindent The off-axis peak flux at viewing angle $q$ relative to the on-axis value $F_{\rm dec}$ is:
\begin{equation}
\frac{F_p(q)}{F_{\rm dec}} = \begin{cases}
\left(\dfrac{\tilde{q}}{\tilde{q}_{\rm dec}}\right)^{-2a_{\rm II}} & q_{\rm dec}\le q<q_j\;\text{(Phase II)}\\[8pt]
\tilde{q}_{\rm dec}^{-2}\left(\dfrac{\tilde{q}}{\tilde{q}_{\rm dec}}\right)^{-2a_{\rm III}} & q\ge q_j\;\text{(Phase III)}
\end{cases}
\end{equation}
Reference fluxes at the characteristic boundaries (used to select the $q_{\rm Euc}$ branch):
\[
F_j = F_{\rm dec}\,\tilde{q}_{\rm dec}^{\,2a_{\rm II}},
\qquad
F_{\rm nr} = F_{\rm dec}\,\tilde{q}_{\rm dec}^{-2}\!\left(\frac{\tilde{q}_{\rm nr}}{\tilde{q}_{\rm dec}}\right)^{-2a_{\rm III}}.
\]
At the boundary $q=q_j$ ($\tilde{q}=1$), the Phase~II and Phase~III expressions are continuous for PLS~G (since $a_{\rm III}^{\rm G}-1=a_{\rm II}^{\rm G}$) but not for PLS~H.

% -------------------------------------------------------
\section{Survey Model}

\begin{equation}
t_{\rm exp} = \frac{f_{\rm live,eff}\,t_{\rm cad}}{N_{\rm exp}} - t_{\rm OH},\qquad
F_{\rm lim} = F_{\rm ref}\!\left(\frac{t_{\rm exp}}{t_{\rm ref}}\right)^{-\alpha},\qquad
f_\Omega = \frac{N_{\rm exp}\,\Omega_{\rm exp}}{4\pi}
\end{equation}
The effective live fraction $f_{\rm live,eff}$ depends on the cadence regime:
\[
f_{\rm live,eff} = \begin{cases}
f_{\rm live} & t_{\rm cad} = n\cdot 86400\,\text{s},\; n\in\mathbb{N}\quad\text{(discrete, per-day)}\\[4pt]
f_{\rm live}/f_{\rm night} & t_{\rm cad} < 86400\,\text{s}\;\text{and optical-survey mode}\quad\text{(continuous, sub-night)}
\end{cases}
\]
with $f_{\rm night}=t_{\rm night}/86400$. In the continuous regime all exposures land inside the night window of length $t_{\rm night}$, so the within-night live fraction is $f_{\rm live}\cdot 86400/t_{\rm night}=f_{\rm live}/f_{\rm night}$. The total detection rate on the continuous branch is additionally multiplied by $f_{\rm night}$ to account for the nighttime accessibility of GRBs. The construction is physical only when $f_{\rm night}\ge f_{\rm live}$ (so $f_{\rm live,eff}\le 1$); historically this was enforced in the UI by a dynamic lower bound on the $t_{\rm night}$ slider.

\medskip\noindent\textbf{Night-schedule generalisation.} With the two-timescale night schedule of Sec.~\ref{sec:schedule}, optical-survey mode is restricted to the discrete cadences $t_{\rm cad}=n\cdot86400$\,s and governed by the single budget equation~\eqref{eq:budget_unified}, parameterised by the within-night usable fraction $f_{\rm eff}$ ($f_{\rm live}\equiv f_{\rm eff}\,f_{\rm night}$). The continuous (sub-day) branch is removed as redundant --- it was exactly the degenerate schedule $(n{=}1,\ N_v{=}t_{\rm night}/\Delta t_v)$, now expressed properly through $N_v$ and $\Delta t_v$. The $f_{\rm live,eff}$ bookkeeping, the two separate model instances it required, the post-hoc $\times f_{\rm night}$ rate factor, and the $f_{\rm night}\ge f_{\rm live}$ constraint (now the trivial $f_{\rm eff}\le1$) are all retired. Non-optical mode keeps the original continuous-cadence model unchanged, parameterised by the wall-clock $f_{\rm live}$ itself: with no night window, $f_{\rm eff}$ is undefined there, so the two modes expose \emph{separate} budget sliders ($f_{\rm eff}$ optical, $f_{\rm live}$ non-optical) rather than one number silently reinterpreted by the mode toggle.

A strategy point $(N_{\rm exp},t_{\rm cad})$ is in the valid region~A0 when $t_{\rm exp}>0$, $N_{\rm exp}\ge1$, and $N_{\rm exp}\le N_{\rm exp,max}\equiv\Omega_{\rm srv,max}/\Omega_{\rm exp}$.

% -------------------------------------------------------
\section{Detection Rate: Two Horizons and Seven Regimes}

\subsection{Distance and horizon angles}

\noindent The maximum detection distance for an on-axis burst (q=0):
\begin{equation}
D_{\rm dec} = D_{\rm euc}\!\left(\frac{F_{\rm lim}}{F_{\rm dec}}\right)^{-1/2}
\end{equation}

\noindent\textbf{Euclidean horizon angle} $q_{\rm Euc}$: the largest off-axis angle detectable at the calibration distance $D_{\rm euc}$:
\begin{equation}
\tilde{q}_{\rm Euc} = \begin{cases}
\tilde{q}_{\rm dec}\!\left(\dfrac{F_{\rm lim}}{F_{\rm dec}}\right)^{-1/(2a_{\rm II})} & F_{\rm lim}\ge F_j\;\text{(Phase II)}\\[8pt]
\tilde{q}_{\rm dec}^{(a_{\rm III}-1)/a_{\rm III}}\!\left(\dfrac{F_{\rm lim}}{F_{\rm dec}}\right)^{-1/(2a_{\rm III})} & F_{\rm lim}<F_j\;\text{(Phase III)}
\end{cases}
\end{equation}
The two branches are continuous at $F_{\rm lim}=F_j$ (for PLS~G), where both give $\tilde{q}_{\rm Euc}=1$ (i.e.\ $q_{\rm Euc}=q_j$).

\noindent\textbf{Off-axis peak time.}
A burst at viewing angle $q$ reaches its afterglow flux peak at time:
\begin{equation}
t_p(q) = \begin{cases}
t_j\,\tilde{q}^{8/3} & q_{\rm dec}\le q<q_j\;\text{(Phase II)} \\[4pt]
t_j\,\tilde{q}^{2}   & q\ge q_j\;\text{(Phase III)}
\end{cases}
\end{equation}
Both expressions equal $t_j$ continuously at $\tilde{q}=1$ ($q=q_j$). At $q=q_{\rm dec}$ ($\tilde{q}=\tilde{q}_{\rm dec}$), Phase~II gives $t_p=t_{\rm dec}$ as expected.

\noindent\textbf{Cadence horizon angle} $q_i$: largest off-axis angle whose peak arrives within $i$ cadence intervals ($t_p=i\,t_{\rm cad}$):
\begin{equation}
\tilde{q}_i = \begin{cases}
(i\,t_{\rm cad}/t_j)^{3/8} & i\,t_{\rm cad}<t_j\;\text{(Phase II peak)}\\[4pt]
(i\,t_{\rm cad}/t_j)^{1/2} & i\,t_{\rm cad}\ge t_j\;\text{(Phase III peak)}
\end{cases}
\end{equation}

\noindent\textbf{Cadence distance} $D_i$: maximum detection distance at angle $q_i$:
\begin{equation}
D_i = \begin{cases}
D_{\rm dec}\,(\tilde{q}_i/\tilde{q}_{\rm dec})^{-a_{\rm II}} & i\,t_{\rm cad}<t_j\\[4pt]
(D_{\rm dec}/\tilde{q}_{\rm dec})\,(\tilde{q}_i/\tilde{q}_{\rm dec})^{-a_{\rm III}} & i\,t_{\rm cad}\ge t_j
\end{cases}
\end{equation}

\subsection{Piecewise rate formulas (dominant-term approximation)}

\begin{table}[h]
\centering
\small
\renewcommand{\arraystretch}{1.4}
\begin{tabular}{cll}
\toprule
Regime & Condition & Rate $R_k$ [yr$^{-1}$] \\
\midrule
A1 & $q_{\rm Euc}\ge q_{\rm nr}$ \textbf{and} $q_{\rm Euc}\ge q_i$ & $f_\Omega\,R_{\rm int}$ \\
A2 & $q_j\le q_{\rm Euc}<q_{\rm nr}$ \textbf{and} $q_{\rm Euc}\ge q_i$ & $\tfrac{1}{2}f_\Omega\,\theta_j^2\,q_{\rm Euc}^2\,R_{\rm int}$ \\
A3 & $q_{\rm dec}\le q_{\rm Euc}<q_j$ \textbf{and} $q_{\rm Euc}\ge q_i$ & $\tfrac{1}{2}f_\Omega\,\theta_j^2\,q_{\rm Euc}^2\,R_{\rm int}$ \\
A4 & $q_i\ge q_{\rm nr}$ \textbf{and} $q_{\rm Euc}<q_i$ & $\tfrac{1}{2}f_\Omega\,\theta_j^2\,q_{\rm nr}^2\,(D_i/D_{\rm euc})^3\,R_{\rm int}$ \\
A5 & $q_j\le q_i<q_{\rm nr}$ \textbf{and} $q_{\rm Euc}<q_i$ & $\tfrac{1}{2}f_\Omega\,\theta_j^2\,q_i^2\,(D_i/D_{\rm euc})^3\,R_{\rm int}$ \\
A6 & $q_{\rm dec}\le q_i<q_j$ \textbf{and} $q_{\rm Euc}<q_i$ & $\tfrac{1}{2}f_\Omega\,\theta_j^2\,q_i^2\,(D_i/D_{\rm euc})^3\,R_{\rm int}$ \\
A7 & $q_{\rm Euc}<q_{\rm dec}$ \textbf{and} $q_i<q_{\rm dec}$ & $\tfrac{1}{2}f_\Omega\,\theta_j^2\,q_{\rm dec}^2\,(D_{\rm dec}/D_{\rm euc})^3\,R_{\rm int}$ \\
\bottomrule
\end{tabular}
\end{table}

\noindent Regimes A2 and A3 share the same formula (both flux-limited); the distinction is only which afterglow phase $q_{\rm Euc}$ falls in. Similarly A5 and A6 share the same formula (both cadence-limited). The general structure is:
\[
R_k = \tfrac{1}{2}f_\Omega\,\theta_j^2\,q_{\rm max}^2\,(D_{\rm eff}/D_{\rm euc})^3\,R_{\rm int},
\]
where $q_{\rm max}$ is the relevant upper angular limit and $D_{\rm eff}$ is the corresponding detection distance. The output surface shows $\log_{10}R_{\rm det}$.

\noindent Region~A0 ($t_{\rm exp}\le0$, $N_{\rm exp}<1$, or $N_{\rm exp}>N_{\rm exp,max}$) is undefined (NaN).
Boundary comparisons use tolerance $\varepsilon=10^{-12}$.

% -------------------------------------------------------
\section{Full-Integral Mode (Optional)}

Instead of the dominant-term formulas, the rate is computed by numerically integrating over all viewing angles using normalized (dimensionless) distances $\tilde{D}=D/D_{\rm euc}$:
\begin{equation}
\label{eq:full_integral}
R = f_\Omega\,\theta_j^2\,R_{\rm int}\int_0^{q_{\rm nr}} q\,\bigl[\min\!\bigl(\tilde{D}_{\rm max}(q),\,\tilde{D}_{\rm eff}\bigr)\bigr]^3\,dq,
\end{equation}
where $\tilde{D}_{\rm eff}=\min(D_i/D_{\rm euc},\,1)$ encodes the cadence constraint and the Euclidean horizon (under the \texttt{win\_from\_peak} setting of Sec.~\ref{sec:window_settings} it becomes the $q$-dependent $\min(\tilde D_w(q),1)$ of Eq.~\eqref{eq:D_window}), and
\begin{equation}
\label{eq:Dmax_tilde}
\tilde{D}_{\rm max}(q) = \begin{cases}
\tilde{D}_{\rm dec} & q<q_{\rm dec}\\[2pt]
\tilde{D}_{\rm dec}(\tilde{q}/\tilde{q}_{\rm dec})^{-a_{\rm II}} & q_{\rm dec}\le q<q_j\\[2pt]
(\tilde{D}_{\rm dec}/\tilde{q}_{\rm dec})(\tilde{q}/\tilde{q}_{\rm dec})^{-a_{\rm III}} & q\ge q_j
\end{cases}
\end{equation}
with $\tilde{D}_{\rm dec}=(F_{\rm lim}/F_{\rm dec})^{-1/2}$ (dimensionless). Both integrands are dimensionless, so~\eqref{eq:full_integral} has units of $R_{\rm int}$ [yr$^{-1}$] as required.

The dominant-term approximation (Sec.~4.2) is the leading term of~\eqref{eq:full_integral} evaluated at the boundary angle ($q_{\rm Euc}$ or $q_i$). The full integral adds contributions from all other angles.
Evaluated with a trapezoidal rule over $N_q=500$ points ($<0.3\%$ error for typical parameters).

% -------------------------------------------------------
\section{Detection-Rate Filters: $q_{\rm min}$ and $D_{\rm min}$}
\label{sec:filters}

The model exposes two physics filters that restrict which GRBs contribute to the detection rate.
They replace an earlier ``off-axis only'' boolean toggle (deprecated; see Sec.~\ref{sec:filters_legacy}).

\medskip\noindent\textbf{Definitions.}
\begin{itemize}
\item $q_{\rm min}\in[0,q_{\rm nr}]$: minimum viewing-angle parameter. Only GRBs with $q\ge q_{\rm min}$ count toward the rate. Default $q_{\rm min}=0$ (no filter).
\item $D_{\rm min}\in[0,D_{\rm euc}]$: minimum \emph{detectable} distance. The integration shell is restricted to $D\ge D_{\rm min}$. Default $D_{\rm min}=0$ (no filter).
\end{itemize}

\medskip\noindent\textbf{Interpretation note.} $D_{\rm min}$ is a \emph{detectability-distance} floor, not a source-distance floor. The detector volume integral is performed over the shell $[\,D_{\rm min},\,D_{\rm eff}(q)\,]$, where $D_{\rm eff}(q)$ is the maximum distance at which a burst at angle $q$ is detectable.
Bursts whose maximum detectable distance $D_{\rm eff}(q)$ falls below $D_{\rm min}$ are excluded.

\subsection{Joint distribution of detected bursts}

For uniform GRB number density and isotropic jet axes, the joint detection PDF on the support
$\{q\in[0,q_{\rm nr}],\;0\le D\le D_{\rm eff}(q)\}$ is
\[
p(q,D)\propto q\cdot D^2,
\]
combining the solid-angle factor $\theta_j^2\,q\,dq$ and the Euclidean volume element $D^2\,dD$.
The unfiltered total rate is
\[
R = f_\Omega\,\theta_j^2\,R_{\rm int}\int_0^{q_{\rm nr}} q\,\tfrac{1}{3}\,D_{\rm eff}(q)^3\,dq,
\]
in agreement with Eq.~\eqref{eq:full_integral} after absorbing the constant $\tfrac{1}{3}$ into the rate normalization.

\subsection{Filtered dominant-term formula}

When $D_{\rm eff}(q)$ is constant over the integration range $[q_{\rm min},q_{\rm max,\,k}]$ (which holds by construction in each regime $k$), the inner $D$-integral becomes
\[
\int_{D_{\rm min}}^{D_{\rm eff,\,k}} D^2\,dD = \tfrac{1}{3}\bigl[D_{\rm eff,\,k}^3 - D_{\rm min}^3\bigr]_+,
\]
where $[\,x\,]_+\equiv\max(x,0)$ enforces non-negativity when the filter excludes the entire shell.
The $q$-integral yields the same algebraic form:
\[
\int_{q_{\rm min}}^{q_{\rm max,\,k}} q\,dq = \tfrac{1}{2}\bigl[q_{\rm max,\,k}^2 - q_{\rm min}^2\bigr]_+.
\]
Combining, the filtered regime rate is
\begin{equation}
\label{eq:filter_dominant}
R_k(q_{\rm min},D_{\rm min}) = \tfrac{1}{2}\,f_\Omega\,\theta_j^2\,R_{\rm int}\;\bigl[q_{\rm max,\,k}^2 - q_{\rm min}^2\bigr]_+\;\bigl[(D_{\rm eff,\,k}/D_{\rm euc})^3 - (D_{\rm min}/D_{\rm euc})^3\bigr]_+.
\end{equation}
At $q_{\rm min}=D_{\rm min}=0$ this reduces exactly (to floating-point) to the unfiltered formula in Sec.~4.2. When either bracket evaluates to zero, the regime contributes nothing and the cell's $\log_{10}R$ is NaN.

\subsection{Filtered full-integral formula}

In full-integral mode the inner $D$-integral cannot be evaluated in closed form because $D_{\rm eff}(q)$ is piecewise:
\begin{equation}
\label{eq:filter_full}
R = f_\Omega\,\theta_j^2\,R_{\rm int}\int_0^{q_{\rm nr}} \mathbf{1}\{q\ge q_{\rm min}\}\;q\;\bigl[D_{\rm eff}(q)^3 - (D_{\rm min}/D_{\rm euc})^3\bigr]_+\,dq,
\end{equation}
where $D_{\rm eff}(q)=\min\!\bigl(\tilde D_{\rm max}(q),\,\tilde D_{\rm eff}\bigr)$ is the same profile defined in Eq.~\eqref{eq:Dmax_tilde}.
The indicator $\mathbf{1}\{q\ge q_{\rm min}\}$ is applied to a discrete $q$-grid of $N_q=500$ points; the trapezoidal rule near $q=q_{\rm min}$ introduces an $\mathcal{O}(\Delta q/q_{\rm nr})\approx 2\times 10^{-3}$ relative bias. This bias is verified empirically by the test \texttt{test\_B12\_indicator\_bias\_at\_qmin\_bounded}.

\subsection{Filtered medians}

Under the dominant-term assumption (constant $D_{\rm eff}$ on the regime), the marginals are
\[
p(q)\propto q\cdot\bigl[D_{\rm eff}^3-D_{\rm min}^3\bigr]_+\propto q,\quad q\in[q_{\rm min},q_{\rm max}],
\]
\[
p(D)\propto D^2\cdot\bigl[q_{\rm max}^2-q_{\rm min}^2\bigr]_+\propto D^2,\quad D\in[D_{\rm min},D_{\rm eff}].
\]
Solving $F(q_{\rm med})=\tfrac{1}{2}$ for each gives
\begin{equation}
\label{eq:filter_medians}
q_{\rm med} = \sqrt{\tfrac{q_{\rm max}^2 + q_{\rm min}^2}{2}},
\qquad
D_{\rm med} = \!\left(\tfrac{D_{\rm eff}^3 + D_{\rm min}^3}{2}\right)^{\!1/3}\!.
\end{equation}
At $q_{\rm min}=D_{\rm min}=0$ these reduce to $q_{\rm max}/\sqrt{2}$ and $D_{\rm eff}\cdot 2^{-1/3}$ (Sec.~6.1). When the filter excludes all weight ($q_{\rm max}\le q_{\rm min}$ or $D_{\rm eff}\le D_{\rm min}$), both medians return NaN.

In full-integral mode where $D_{\rm eff}(q)$ varies, the q-marginal becomes
\[
p(q)\propto q\cdot\bigl[D_{\rm eff}(q)^3-(D_{\rm min}/D_{\rm euc})^3\bigr]_+\cdot \mathbf{1}\{q\ge q_{\rm min}\},
\]
which is \emph{not} simply $\propto q$, so the analytic formula~\eqref{eq:filter_medians} no longer applies; the median is found by cumulative trapezoidal integration of this weighted form. See \texttt{compute\_medians\_numerical}.

\subsection{Fading-rate filter $s_{\rm fade}$}
\label{sec:filter_sfade}

Real optical transient surveys filter their alert streams to ``afterglow-like'' candidates by requiring a measured magnitude decline faster than some threshold $s_{\rm fade}$ in mag/day; common values are $0.3$ (ZTF-style) and $1$ (stricter). The filter is exposed as a third sidebar parameter alongside $q_{\rm min}$ and $D_{\rm min}$, with two modes that coincide in the $t_{\rm cad}\to 0$ limit.

\paragraph{Setup.} Post-peak, the lightcurve obeys $F(t)\propto t^\alpha$ with $\alpha<0$, so
\[
m(t)=-2.5\alpha\log_{10}(t/t_p)+\text{const},\qquad
\frac{dm}{dt}=-\frac{2.5\,\alpha}{t\,\ln 10}\;\text{mag/s}.
\]
The survey samples this lightcurve starting at the first post-peak detection $t_{\rm first}\ge t_p$ and ending at $t_{\rm last}=t_{\rm first}+(i-1)\,t_{\rm cad}$. A key identity: for a pure power law the measured slope between $t_{\rm first}$ and $t_{\rm last}$ depends only on $t_{\rm first}$ and the baseline $(i-1)\,t_{\rm cad}$ --- not on $t_p$ separately --- so the cut inverts cleanly to a cap on the \emph{first-detection time}.

\paragraph{Discrete (per-cadence) mode.} The measured slope across the $i$ detections is
\begin{equation}
\label{eq:s_meas_discrete}
s_{\rm meas}=-2.5\,\alpha\,\log_{10}\!\left(1+\frac{(i-1)\,t_{\rm cad}}{t_{\rm first}}\right)\frac{\text{DAY\_S}}{(i-1)\,t_{\rm cad}}\;\;\text{[mag/day]}.
\end{equation}
Imposing $s_{\rm meas}\ge s_{\rm fade}$ and solving for the latest $t_{\rm first}$ that still passes, with $A=2.5|\alpha|$ and $E=s_{\rm fade}(i-1)t_{\rm cad}/(\text{DAY\_S}\cdot A)$:
\begin{equation}
\label{eq:tfs_discrete}
t_{f,s}=\frac{(i-1)\,t_{\rm cad}}{10^E-1},\qquad\text{pass}\iff t_{\rm first}\le t_{f,s}.
\end{equation}
This is strategy-dependent: long $t_{\rm cad}$ shrinks $t_{f,s}$ (the survey can no longer resolve the slope). At $i=1$ the slope is undefined and the cut is bypassed.

\paragraph{Continuous (instantaneous) mode.} Evaluate $|dm/dt|$ at the first detection directly and require it to exceed $s_{\rm fade}$:
\begin{equation}
\label{eq:tfs_continuous}
t_{f,s}=\frac{2.5\,|\alpha|\,\text{DAY\_S}}{s_{\rm fade}\,\ln 10}.
\end{equation}
Strategy-independent; defined for all $i$.

\paragraph{Limit equivalence.} As $t_{\rm cad}\to 0$, $E\to 0$ and $10^E-1\approx E\ln 10$, so~\eqref{eq:tfs_discrete} reduces to~\eqref{eq:tfs_continuous}:
\[
t_{f,s}\to\frac{(i-1)\,t_{\rm cad}}{E\ln 10}=\frac{2.5\,|\alpha|\,\text{DAY\_S}}{s_{\rm fade}\,\ln 10}.
\]
The continuous mode is the high-cadence limit of the discrete mode.

\paragraph{Dominant-term treatment: best-case start.} In the closed-form (dominant-term) rate mode we assume the first detection lands at the peak, $t_{\rm first}\approx t_p(\tilde q)$, so $t_{f,s}$ caps the peak time itself. The model's $q_i$ relation (\S\,Viewing-angle parameterisation) gives $t_p(\tilde q)=t_j\tilde q^{8/3}$ in Phase~II and $t_j\tilde q^2$ in Phase~III. Inverting:
\begin{equation}
\label{eq:qs_phase}
\tilde q_{s,\rm II}=\left(\frac{t_{f,s,\rm II}}{t_j}\right)^{3/8},\qquad
\tilde q_{s,\rm III}=\left(\frac{t_{f,s,\rm III}}{t_j}\right)^{1/2},
\end{equation}
where $\alpha$ in $t_{f,s}$ is the phase-specific temporal slope: $\alpha_{\rm II}=-3(p-1)/4$, $\alpha_{\rm III}=-p$ for PLS~G, and $\alpha_{\rm II}=-(3p-2)/4$, $\alpha_{\rm III}=-p$ for PLS~H. The cut then acts on each regime's effective $q_{\rm max}$ as a phase-mapped hard cap: the seven dominant-term rates~\eqref{eq:filter_dominant} become
\[
r_k = \tfrac{1}{2}f_\Omega\theta_j^2 R_{\rm int}\cdot
\bigl[\min(q_{\rm max,\,k},\,q_{s,\,\Phi(k)})^2-q_{\rm min}^2\bigr]_+\cdot
\bigl[(D_{\rm eff,\,k}/D_{\rm euc})^3-(D_{\rm min}/D_{\rm euc})^3\bigr]_+,
\]
with phase assignment $\Phi$: A1, A2, A4, A5 $\to$ Phase~III; A3, A6, A7 $\to$ Phase~II. As shown next, this best-case treatment is the exact $t_{\rm cad}\to 0$ limit of the uniform-start model and an optimistic upper bound otherwise --- consistent with the dominant-term philosophy of keeping the rate expressions closed-form.

\paragraph{Uniform observation start (full-integral mode).} The survey's visit schedule has a random phase relative to the burst, so the first post-peak visit actually lands at
\[
t_{\rm first}=t_p(\tilde q)+u,\qquad u\sim\mathcal U(0,\,t_{\rm cad}).
\]
Since pass $\iff t_{\rm first}\le t_{f,s}$, each burst passes the cut with probability
\begin{equation}
\label{eq:P_fade}
P(q)=\operatorname{clip}\!\left(\frac{t_{f,s}-t_{p,\rm eff}(q)}{t_{\rm cad}},\,0,\,1\right),\qquad
t_{p,\rm eff}(q)=\max\!\bigl(t_j\,[\max(q-1,0)]^{\beta},\;t_{\rm dec}\bigr),
\end{equation}
with $\beta=8/3$ (Phase~II) and $\beta=2$ (Phase~III). $P$ reaches $0$ exactly at the hard cap $q_s$ of~\eqref{eq:qs_phase} and ramps up (linearly in $t_p$) below it. The $t_{\rm dec}$ floor on the effective peak time is \emph{exact} for on-axis viewers --- $t_j\,\tilde q_{\rm dec}^{8/3}=t_{\rm dec}$ identically, so the floor binds precisely for $q\le q_{\rm dec}$ (where the model lightcurve peaks at $t_{\rm dec}$) and is a no-op above; on-axis viewers are lumped with Phase~II as in the existing masks. For the fade weight the floor matters only at $\mathcal O(t_{\rm dec}/t_{\rm cad})\approx 19\,\mathrm s/t_{\rm cad}$ (it is however \emph{essential} for the rise filter, Sec.~\ref{sec:filter_srise}); one physical consequence: when $t_{f,s}<t_{\rm dec}$ --- i.e.\ even a burst caught at the earliest possible epoch fades too slowly --- the weight now correctly annihilates the rate instead of leaving the fictitious pre-$t_{\rm dec}$ residual of the pure extrapolation. In the full-integral form~\eqref{eq:filter_full} the integrand simply picks up the factor $P(q)$ in place of the former indicator $\mathbf 1\{q\le q_{s,\Phi(q)}\}$; this weight is applied only where the integral is evaluated numerically anyway (\texttt{rate\_log10\_full\_integral}, the numerical medians, and the $dR/dq$, $dR/dD$ views), leaving the closed-form rate equations untouched.

Note the qualitatively new effect: when $t_{f,s}<t_{\rm cad}$ even on-axis bursts pass only with probability $t_{f,s}/t_{\rm cad}$ --- with a long cadence the afterglow is usually caught late, when its mag/day decline is slow. The best-case hard cap would let all of them pass.

\paragraph{Limits and invariants.}
(i)~$s_{\rm fade}\to 0$: $t_{f,s}=+\infty$, $P\equiv 1$, and both rate modes collapse to the unfiltered baseline to floating point.
(ii)~$t_{\rm cad}\to 0$: the ramp width vanishes and $P\to\mathbf 1\{q\le q_s\}$, recovering the dominant-term hard cap; the deficit is linear in $t_{\rm cad}$.
(iii)~$P(q)\le\mathbf 1\{q\le q_s\}$ pointwise, so under an active cut the full-integral rate sits \emph{below} the dominant-term rate --- an expected, physical mode disagreement, not a numerical artefact.
(iv)~$t_{f,s}$ is strictly decreasing in $s_{\rm fade}$, so both treatments remain monotone non-increasing in $s_{\rm fade}$.

\paragraph{Assumptions.} (a)~The start phase $u$ is treated as independent of the detection condition itself; the $q_i$ detection machinery keeps its dominant-term form. This is coherent: $u\le t_{\rm cad}$ implies $t_{\rm last}\le t_p+i\,t_{\rm cad}$, matching the $q_i$ detection window. (b)~The rise before the peak is ignored --- the slope is measured from the first post-peak detection. (c)~The measurement segment is treated as lying in a single power-law phase; a delayed segment that straddles $t_j$ actually fades \emph{faster} than modeled (the post-break slope is steeper), so the treatment is conservative.

\paragraph{Implementation reference.} \texttt{DetectionRateModel.\_t\_first\_caps()} returns the phase-resolved $t_{f,s,\rm II}$ and $t_{f,s,\rm III}$ arrays (broadcastable to the strategy grid). \texttt{\_q\_s\_fading\_caps()} maps them to the best-case hard caps~\eqref{eq:qs_phase} used by \texttt{rate\_log10} and \texttt{compute\_medians\_analytic}; \texttt{\_fading\_survival()} evaluates the pointwise $P(q)$ of~\eqref{eq:P_fade} used by \texttt{rate\_log10\_full\_integral}, \texttt{compute\_medians\_numerical}, \texttt{dR\_dq\_full\_integral}, and \texttt{dR\_dD\_full\_integral}. Since $P$ is a continuous weight rather than a cut, the numerical median integrals remain valid without any non-contiguous-interval special-casing.

\subsection{Rise-rate filter $s_{\rm rise}$}
\label{sec:filter_srise}

Alert streams also cut on the \emph{rise}: a new candidate is credible when the previous visit --- one cadence before the first detection --- was a non-detection, bounding the earlier magnitude by the limiting magnitude $m_{\rm lim}$. The implied rise rate is then at least $(m_{\rm lim}-m_{\rm first})/t_{\rm cad}$, and the filter requires it to exceed $s_{\rm rise}$ [mag/day].

\paragraph{Semantics.} With $m=-2.5\log_{10}F+\text{const}$,
\begin{equation}
\label{eq:eta_rise}
\frac{m_{\rm lim}-m_{\rm first}}{t_{\rm cad}/\text{DAY\_S}}\ge s_{\rm rise}
\iff
F(t_{\rm first})\ge\eta\,F_{\rm lim},\qquad
\eta=10^{E},\quad E=\frac{s_{\rm rise}\,t_{\rm cad}}{2.5\,\text{DAY\_S}}.
\end{equation}
The boost $\eta$ grows \emph{exponentially} with cadence: a slow survey can only certify a fast riser if the source jumped far above the limit since the last visit. The model ignores the pre-peak rise, so the previous visit (always pre-peak, $t_{\rm first}-t_{\rm cad}<t_p$) is automatically a non-detection --- the filter is internally consistent with the baseline. Unlike the discrete fade filter, the rise cut applies at $i_{\rm det}=1$ (it needs only the previous visit and the first detection); it has no discrete/continuous mode split --- the measurement \emph{is} a previous-visit finite difference, so \texttt{s\_mode} remains fade-only.

\paragraph{Dominant-term treatment: best-case start.} With the first detection at the peak, pass $\iff F_p(q,D)\ge\eta F_{\rm lim}\iff D\le D_{\rm max}(q)/\sqrt\eta$. Because the flux-limited boundary functions take $F_{\rm lim}$ as an argument, the raised-limit boundaries are direct evaluations:
\[
q_{E,r}=q_{\rm Euc}(\eta F_{\rm lim}),\qquad
q_{r,i}=q_{\rm Euc}\!\bigl(\eta F_{\rm lim}\,(D_i/D_{\rm euc})^2\bigr),\qquad
D_{{\rm dec},r}=D_{\rm dec}\,\eta^{-1/2},
\]
where $q_{r,i}$ is the angle at which $D_{\rm max}(q)/\sqrt\eta=D_i$; since $D_i=D_{\rm max}(q_i)$ identically, $q_{r,i}(\eta{=}1)=q_i$ exactly. The caps enter the regime formulas as $q_{\rm max}\to\min(q_{\rm max},\,q_{E,r})$ for the flux-limited A1--A3, $\min(q_{\rm max},\,q_{r,i})$ for the cadence-limited A4--A6, and $D_{\rm eff,7}\to\min(D_{\rm dec},D_i,D_{{\rm dec},r})$ for A7, composing with the fading caps of~\eqref{eq:qs_phase} by further $\min()$.

\paragraph{Deep-cut correctness (on-axis max-term).} $q_{\rm Euc}$ extrapolates the off-axis power law and never returns $q<1$, so for a deep cut the truncated angular term would saturate at $[1-q_{\rm min}^2]_+\times$D-factor instead of vanishing. Moreover the extrapolated boundary is the \emph{true} flux boundary only while it sits above $q_{\rm dec}$: below the crossing, the on-axis $D_{\rm max}=D_{\rm dec}$ plateau makes the angular box overcount. Each of $r_1$--$r_6$ is therefore computed as the \textbf{max of two closed-form lower bounds} of the true pass volume:
\[
r_k=\text{base}\cdot\max\Bigl(
\underbrace{\mathbf 1\{q_{{\rm cap},k}\ge q_{\rm dec}\}\,\bigl[q_{{\rm max},k,\rm eff}^2-q_{\rm min}^2\bigr]_+\,\bigl[\tilde D_{{\rm eff},k}^3-\tilde D_{\rm min}^3\bigr]_+}_{\text{angular term (valid above the crossing)}},\;
\underbrace{\bigl[\min(q_{\rm dec},q_{{\rm max},k},q_{s,\rm II})^2-q_{\rm min}^2\bigr]_+\,\bigl[\min(\tilde D_{{\rm eff},k},\tilde D_{\rm dec}\,\eta^{-1/2})^3-\tilde D_{\rm min}^3\bigr]_+}_{\text{on-axis term}}
\Bigr).
\]
At the crossing $q_{{\rm cap},k}=q_{\rm dec}$ one has $D_{\rm dec}/\sqrt\eta=D_{\rm eff}$-scale, so the two terms are equal and the max is continuous in $s_{\rm rise}$; at $\eta=1$ the on-axis term is dominated and the baseline is recovered; in the deep-cut limit $r_k\propto\eta^{-3/2}\to 0$ (verified numerically over 6~dex).

\paragraph{Uniform observation start (full-integral mode).} Post-peak, $F(t)=F_p\,(t/t_p)^{-|\alpha|}$ crosses the raised limit at
\begin{equation}
\label{eq:t_lim}
t_{\rm lim}(q,\tilde D)=t_{p,\rm eff}(q)\left(\frac{\tilde D_0(q)}{\tilde D}\right)^{k},\qquad
k=\frac{2}{|\alpha_{\Phi}|},\qquad
\tilde D_0(q)=\tilde D_{\rm max}(q)\,\eta^{-1/2},
\end{equation}
so pass $\iff t_{\rm first}\le t_{\rm lim}$. Here the $t_{\rm dec}$ floor of $t_{p,\rm eff}$~\eqref{eq:P_fade} is \emph{essential}: $t_{\rm lim}\propto t_{p,\rm eff}$, and the pure extrapolation ($t_p\to 0$ on-axis) would wrongly annihilate all on-axis rise weights; with the floor, the general formula reproduces the on-axis lightcurve $F=F_{\rm dec}(D)\,(t/t_{\rm dec})^{-|\alpha_{\rm II}|}$ exactly.

\textbf{Joint weight --- same start offset.} Both cuts are conditions on the \emph{same} uniform offset $u$ (comonotone, not independent; a product would double-penalize):
\begin{equation}
\label{eq:w_joint}
w(q,\tilde D)=\operatorname{clip}\!\left(\frac{\min\bigl(t_{f,s,\Phi}(q),\,t_{\rm lim}(q,\tilde D)\bigr)-t_{p,\rm eff}(q)}{t_{\rm cad}},\,0,\,1\right).
\end{equation}
Because $w$ now depends on $\tilde D$, the $q$-integrand's distance volume is no longer $[\tilde D_{\rm eff}^3-\tilde D_{\rm min}^3]_+\,P$; it becomes the \textbf{closed-form weighted volume}
\begin{equation}
\label{eq:V_w}
V_w(q)=\int_{\tilde D_{\rm min}}^{\tilde D_{\rm eff}(q)}3\tilde D^2\,w(q,\tilde D)\,d\tilde D
=\frac{w_f}{t_{\rm cad}}\Bigl[\min(\tilde D_{\rm eff},\tilde D_c)^3-\tilde D_{\rm min}^3\Bigr]_+
+\frac{t_{p,\rm eff}}{t_{\rm cad}}\bigl[G(U)-G(L)\bigr]_+,
\end{equation}
with the fade window $w_f=\operatorname{clip}(t_{f,s}-t_{p,\rm eff},0,t_{\rm cad})$, ramp start $\tilde D_c=\tilde D_0\,[t_{p,\rm eff}/(t_{p,\rm eff}+w_f)]^{1/k}$, antiderivative $G(\tilde D)=3\tilde D_0^{k}\tilde D^{3-k}/(3-k)-\tilde D^3$ (of $3\tilde D^2[(\tilde D_0/\tilde D)^k-1]$), and clamped ramp edges $U=\min(\tilde D_{\rm eff},\tilde D_0)$, $L=\min(\max(\tilde D_{\rm min},\tilde D_c),U)$. The exponent domain is safe: $k_{\rm II}=8/[3(p-1)]\le 2.64$ and $k_{\rm III}=2/p<1$ for the app's $p\in[2.01,3]$, so the $k=3$ pole is unreachable. The pointwise $w$~\eqref{eq:w_joint} drives the $dR/dD$ view and the numerical $D$-median; $V_w$~\eqref{eq:V_w} drives the rate integral, the numerical $q$-median, and the $dR/dq$ view. Numerics: $\eta^{-1/2}$ is computed directly as $10^{-E/2}$ so that extreme $s_{\rm rise}\,t_{\rm cad}$ (where $\eta$ overflows) degrades gracefully to $\tilde D_0=0$, $V_w=0$.

\paragraph{Limits and invariants.}
(i)~$s_{\rm rise}=0$ bypasses everything bit-identically (both modes, all derived quantities).
(ii)~Monotone non-increasing in $s_{\rm rise}$, and jointly in $s_{\rm fade}$.
(iii)~Pointwise $w\le\min(P_{\rm fade},\,\mathbf 1\{\tilde D\le\tilde D_0\})$, so the joint rate sits below each single-cut rate, and the weighted full-integral rate sits below the best-case ($u=0$) hard-cut rate built from the same $\tilde D_{\rm max}$ profile. \emph{Note}: full vs dominant themselves are \textbf{not} ordered under the rise cut --- the dominant truncation drops the $D$-structured tail beyond the boundary angle, which can outweigh the uniform-start suppression the full mode applies.
(iv)~$s_{\rm rise}\to 0^{+}$ does \emph{not} limit to the bypass near the flux boundary ($\tilde D\to\tilde D_{\rm max}$): any $s_{\rm rise}>0$ also switches on the real ``catch it before it fades to the limit'' physics that the baseline neglects (sources near their maximum detectable distance are above the limit only briefly). $s_{\rm rise}=0$ disables previous-visit modeling entirely; the residual is confined to the near-boundary shell.

\paragraph{Assumptions.} (a)~The rise itself is ignored (dominant-term lightcurve), which is what makes the previous-visit non-detection automatic. (b)~The start offset $u$ is decoupled from the detection-count machinery ($q_i$, $D_i$ keep their dominant-term forms). (c)~$t_{\rm lim}$ uses the single power-law phase of the viewer; a crossing into the post-$t_j$ (steeper) decay makes the true window slightly \emph{smaller} than modeled, so the rise treatment is mildly permissive there --- opposite in sign to the fade filter's conservative single-phase bias.

\paragraph{Implementation reference.} \texttt{\_rise\_eta()} returns $(\eta,\eta^{-1/2})$; \texttt{\_t\_p\_eff()} the floored peak time; \texttt{\_rise\_fade\_windows()} the shared per-point quantities; \texttt{\_joint\_survival()} the pointwise $w$~\eqref{eq:w_joint}; \texttt{\_weighted\_D\_volume()} the closed form~\eqref{eq:V_w}. The dominant-mode boundary shifts and the on-axis max-term live in \texttt{rate\_log10} (and, in capped-$q_{\rm max}$ form, \texttt{compute\_medians\_analytic}); the bridge's dominant $R(q)$/$R(D)$ view applies the per-regime caps but keeps the separable regime form (no max-term) --- use full-integral mode for the deep-cut corner there.

\subsection{Legacy off-axis-only mode (deprecated)}
\label{sec:filters_legacy}

Earlier versions of the code exposed an ``off-axis detections only'' boolean that subtracted the on-axis ($q<q_{\rm dec}$) contribution: $R_{\rm off}=R_{\rm total}-R_{\rm on}$.
This was replaced by the continuous $q_{\rm min}$ filter described above.
Setting $q_{\rm min}\approx q_{\rm dec}\approx 1.03$ approximates the old behaviour: the angular shell $q<q_{\rm dec}$ is excluded directly from each regime's integral, rather than computed and subtracted.
The new formulation is numerically cleaner (no near-cancellation of two large numbers) and exposes the filter as a continuous physics parameter rather than a hard toggle.

% -------------------------------------------------------
\section{Detection-Window Settings}
\label{sec:window_settings}

Two independent settings refine how the ``$\ge i$ detections within the cadence'' requirement is counted. Both on gives the corrected criterion
\begin{equation}
\label{eq:window_corrected}
t_{+}(D)-t_{p,\rm eff}(q)\;\ge\;(i-1)\,t_{\rm cad}.
\end{equation}
\textbf{Defaults.} The app defaults to \texttt{win\_i\_minus\_one} ON (the $(i-1)$ convention below, in both survey modes); \texttt{win\_from\_peak} is off by default and belongs to the app's Exact rate mode. The engine constructor's own defaults keep both off (the legacy simple mode, bit-for-bit parity-tested) so that library callers and archived baselines are unaffected by the app-level default.

\subsection{Setting 1: $(i-1)$ cadence intervals (\texttt{win\_i\_minus\_one})}

$i$ detections spaced $t_{\rm cad}$ apart bracket only $i-1$ gaps, so the required detectability duration is
\[
T_{\rm req}=(i-1)\,t_{\rm cad}
\qquad\text{instead of}\qquad
T_{\rm req}=i\,t_{\rm cad},
\]
substituted everywhere $T_{\rm req}$ enters ($q_i$, $D_i$, and the window-mode $D_{\rm eff}$ below; helper \texttt{\_t\_req\_s()}). Exactly equivalent to the legacy criterion evaluated at $i-1$. Degenerate case $i_{\rm det}=1$: $T_{\rm req}=0$, so $q_i\to1$ and $D_i\to+\infty$ --- a single detection at the peak needs no duration and the cadence constraint vanishes.

At $N_v\ge2$ the toggle survives only where a hard endpoint window is still used --- the per-channel rectangles of the exact-rectangle mode (span $S_c$ vs $S_c+g_c$, Sec.~\ref{sec:sched_exact}); the gap-averaged dominant mixture and the from-peak phase ramp \emph{integrate} the waiting time instead of bracketing it, so the convention choice dissolves there (Sec.~\ref{sec:sched_dominant}).

\subsection{Setting 2: window measured from the peak (\texttt{win\_from\_peak})}

The legacy criterion measures the window from the \emph{burst}: $t_{+}(D)\ge T_{\rm req}$, with the boundary pair $q_i$ ($t_p=T_{\rm req}$) and $D_i$ ($t_+=T_{\rm req}$) forming the admissible rectangle $\{q\le q_i\}\times\{D\le D_i\}$. But the off-axis flux is negligible before the peak, so epochs in $[0,t_p]$ cannot be detections: at the dominant corner $(q_i,D_i)$ the true above-threshold window $t_+-t_p$ is $\approx 0$ even though $i$ detections were credited. The corrected requirement counts the window from the turn-on $t_{\rm on}\approx t_p$:
\begin{equation}
\label{eq:window_from_peak}
t_{+}(D)\;\ge\;t_{p,\rm eff}(q)+T_{\rm req},
\end{equation}
which turns the cadence cap into the $q$-dependent distance
\begin{equation}
\label{eq:D_window}
D_w(q)=D\bigl(t_+{=}\,t_{p,\rm eff}(q)+T_{\rm req}\bigr),\qquad
D(t_+)=
\begin{cases}
D_{\rm dec}\,(t_+/t_{\rm dec})^{\alpha_{\rm II}/2}, & t_+<t_j,\\[4pt]
D_{\rm dec}\,(t_j/t_{\rm dec})^{\alpha_{\rm II}/2}\,(t_+/t_j)^{\alpha_{\rm III}/2}, & t_+\ge t_j,
\end{cases}
\end{equation}
with $\alpha_\Phi<0$ the temporal flux indices (\texttt{D\_from\_t\_plus()}, \texttt{\_D\_eff\_window\_cm()}). The inversion $D(t_+)$ unifies the two legacy rectangle boundaries: at $t_+=T_{\rm req}$ it reproduces $D_i$ exactly, and at $t_+=t_p(q)$ it reproduces the flux-limited $D_{\rm max}(q)$. Since $t_{p,\rm eff}+T_{\rm req}>t_p$, one has $D_w(q)<D_{\rm max}(q)$ everywhere --- the window cap is always the binding distance constraint, and the full-integral integrand becomes $q\,[\min(\tilde D_w(q),1)]^3$.

\paragraph{Consequences.}
(i)~\textbf{No dominant-term closed form}: $D_{\rm eff}$ is $q$-dependent, and the leading integral $\int q\,\tilde D_w(q)^3\,dq$ is elementary only in phase III. \texttt{rate\_log10} therefore evaluates the $q$-integral internally under this setting (in both UI modes), \texttt{compute\_medians} always takes the numerical path, and the bridge's $R(q)$/$R(D)$ views use the full-integral branch. Region masks keep the legacy $q_i$ boundaries as an approximate classification (regime colouring only).
(ii)~\textbf{Off-axis suppression}: in the phase-III cadence-limited regime the corrected integral evaluates to $1/(3p/2-1)$ of the legacy rectangle ($\approx 0.43$ at $p=2.2$; on-axis contributions are unchanged since $t_{p,\rm eff}=t_{\rm dec}\ll T_{\rm req}$), and the dominant angle shifts inward from $\tilde q_i$ to $\tilde q_i/\sqrt{3p-1}$ --- the detected $q$-median moves on-axis, toward the observed ZTF $q_{\rm med}\lesssim 1$.
(iii)~\textbf{Limits}: $t_{\rm cad}\to0$ gives $D_w(q)\to D_{\rm max}(q)$ and the legacy rate is recovered; the $s_{\rm fade}$/$s_{\rm rise}$ weights compose unchanged (they constrain the first-detection timing, the window constrains the duration).

\paragraph{Implementation reference.} Flags on \texttt{DetectionRateModel} (constructor kwargs \texttt{win\_i\_minus\_one} / \texttt{win\_from\_peak}, threaded through \texttt{make\_rate\_model}; bridge params \texttt{win\_iminus1} / \texttt{win\_tp}; two Settings switches in the app). Tests: \texttt{tests/test\_window\_toggles.py} (parity, the $i\!-\!1$ shift identity, $D(t_+)$ boundary identities, suppression/limit checks).

% -------------------------------------------------------
\section{Median Viewing Angle and Distance}

The model computes the median viewing angle $q_{\rm med}$ and median detection distance $D_{\rm med}$ of the detected burst population at each grid point. Two computation modes are available, matching the rate mode.

\subsection{Analytic medians (dominant-term mode)}

In the dominant-term approximation, the effective detection volume is uniform within an angular cap $[0,q_{\rm max}]$ at constant distance $D_{\rm eff,\,const}$. The marginal distributions are:
\[
p(q)\propto q, \quad q\in[0,\,q_{\rm max}]\qquad\Longrightarrow\qquad q_{\rm med} = \frac{q_{\rm max}}{\sqrt{2}},
\]
\[
p(D)\propto D^2, \quad D\in[0,\,D_{\rm eff,\,const}]\qquad\Longrightarrow\qquad D_{\rm med} = D_{\rm eff,\,const}\cdot 2^{-1/3}.
\]
The quantities $q_{\rm max}$ and $D_{\rm eff,\,const}$ depend on the active regime:

\begin{table}[h]
\centering
\small
\renewcommand{\arraystretch}{1.3}
\begin{tabular}{ccc}
\toprule
Regime & $q_{\rm max}$ & $D_{\rm eff,\,const}$ \\
\midrule
A1 & $q_{\rm nr}$ & $D_{\rm euc}$ \\
A2, A3 & $q_{\rm Euc}$ & $D_{\rm euc}$ \\
A4 & $q_{\rm nr}$ & $D_i$ \\
A5, A6 & $q_i$ & $D_i$ \\
A7 & $q_{\rm dec}$ & $D_{\rm dec}$ \\
\bottomrule
\end{tabular}
\end{table}

\noindent The $p(q)\propto q$ distribution arises from the volume element $\theta_j^2\,q\,dq$ in the solid-angle integral. The $p(D)\propto D^2$ distribution reflects uniform 3-D number density at fixed angular slice.

\subsection{Numerical medians (full-integral mode)}

When full-integral mode is active, medians are computed from the same $\tilde{D}_{\rm eff}(q)$ profile used in Eq.~\eqref{eq:full_integral}.

\medskip\noindent\textbf{Median $q$:}
Define the unnormalized weight $w(q)=q\cdot\tilde{D}_{\rm eff}(q)^3$. Then $q_{\rm med}$ is found from the cumulative integral:
\begin{equation}
\frac{\int_0^{q_{\rm med}} w(q)\,dq}{\int_0^{q_{\rm nr}} w(q)\,dq} = \frac{1}{2}.
\end{equation}
Evaluated by cumulative trapezoidal sum over $N_q=200$ points.

\medskip\noindent\textbf{Median $D$:}
Define $Q(D)=\max\{q:\tilde{D}_{\rm eff}(q)\ge D/D_{\rm euc}\}$ (the largest angle still detectable at distance~$D$). Since $\tilde{D}_{\rm eff}(q)$ is non-increasing, $Q(D)$ is obtained by counting how many grid points satisfy the threshold. The marginal density is:
\begin{equation}
p(D)\propto D^2\,Q(D)^2,
\end{equation}
which reduces to $D^2$ for uniform density and constant $Q$. The median is found by cumulative summation over a grid of 100 normalized distance levels $D/D_{\rm euc,\,eff}\in[0,1]$, where $D_{\rm euc,\,eff}=\tilde{D}_{\rm eff}(0)\cdot D_{\rm euc}$.

% -------------------------------------------------------
\section{Optical Cadence Constraints}
\label{sec:optical_constraints}

For optical surveys, the only physically meaningful cadences are whole numbers of nights:
\[
t_{\rm cad}=n\times1\;\text{day},\qquad n\in\mathbb{Z}^+ ,
\]
enforced as a validity mask (not by rounding); all $t_{\rm cad}<1$\,day grid points are invalid in optical mode. Sub-day sampling is not a separate cadence regime --- revisiting a field within the night is exactly what the schedule parameters $(N_v,\Delta t_v)$ of Sec.~\ref{sec:schedule} describe, on the $n=1$ row. The schedule feasibility masks of Sec.~\ref{sec:sched_feasibility} apply in addition.

\paragraph{Historical note.} Earlier versions modelled a \emph{continuous} sub-night band ($t_{\rm cad}<t_{\rm night}/i$: all $i$ detections in one night, rate multiplied post-hoc by $f_{\rm night}$ via a second model instance \texttt{model\_night}) and declared the band $t_{\rm night}/i\le t_{\rm cad}<1$\,day a forbidden \emph{gap}. Both are subsumed: the continuous band was the degenerate schedule $(n{=}1,\ N_v{=}t_{\rm night}/\Delta t_v)$ with the $f_{\rm night}$ factor now \emph{derived} from the phase geometry (Sec.~\ref{sec:sched_Pi}, short-event limit), and the gap band corresponded to $N_v<i$ on the $n{=}1$ row, where $P_i$ correctly accumulates detections across consecutive nights instead of forbidding the strategy.

% -------------------------------------------------------
\section{Two-Timescale Night Schedule ($N_v$ visits per night)}
\label{sec:schedule}

\subsection{Motivation}

\paragraph{The idea in one paragraph.} A real optical survey visits a field in \emph{bursts of attention}: a couple of images tonight, hours apart, then nothing until the schedule brings the field back some nights later. A transient therefore doesn't need to survive a full inter-night cadence to be caught twice --- if it is up when the survey arrives, the second image lands \emph{hours} later, not days. The one-timescale model of the preceding sections cannot say this: it has a single gap $t_{\rm cad}$ between any two visits, so ``two detections'' always costs two cadences of survival. The night schedule splits the two timescales --- \emph{how often the survey returns to the field} ($t_{\rm cad}=n$ nights) and \emph{how it samples the field within a visit-night} ($N_v$ visits spaced $\Delta t_v$) --- and then re-derives everything the old single number was doing from the actual visit times.

The survey model of the preceding sections assumes that consecutive visits to the same field are always separated by the one cadence $t_{\rm cad}$. That single number then plays five physically distinct roles:
\begin{center}
\small
\renewcommand{\arraystretch}{1.25}
\begin{tabular}{lll}
\toprule
Role & Formula & Where \\
\midrule
exposure budget & $t_{\rm exp}=f_{\rm live,eff}\,t_{\rm cad}/N_{\rm exp}-t_{\rm OH}$ & Survey Model \\
required duration & $T_{\rm req}=i\,t_{\rm cad}$ (or $(i{-}1)\,t_{\rm cad}$) & $q_i$, $D_i$, $D_w$ \\
fade baseline & $(i-1)\,t_{\rm cad}$ & Eq.~\eqref{eq:tfs_discrete} \\
rise gap & $\eta=10^{s_{\rm rise}t_{\rm cad}/(2.5\,\mathrm{DAY})}$ & Eq.~\eqref{eq:eta_rise} \\
visit phase & $u\sim\mathcal U(0,t_{\rm cad})$ & Eqs.~\eqref{eq:P_fade},~\eqref{eq:w_joint} \\
\bottomrule
\end{tabular}
\end{center}
Real optical surveys are \emph{two-timescale}: the ZTF public survey visits each field $N_v=2$ times per night, with intra-night spacing $\Delta t_v\sim0.5$--$5$\,hr, but returns to that field only every $2$ nights; the ZTF high-cadence programs take $6$ visits per night, every night. Neither schedule is expressible in the one-timescale model. Forcing the public mode onto $(t_{\rm cad}=2\,\mathrm d,\,i=2)$ sets $T_{\rm req}=i\,t_{\rm cad}\approx4$\,d for what is physically a pair of detections a few hours apart; since the cadence horizon shrinks as a power of $T_{\rm req}$ ($\tilde q_i\propto T_{\rm req}^{3/8\text{ or }1/2}$, $D_i\propto\tilde q_i^{-a_\Phi}$), the modelled horizon $D_i$ is far too small. This is the leading structural candidate for the too-short predicted median distance of the public mode ($D_{\rm med}\approx1.0$\,Gpc predicted vs.\ $\approx3.7$\,Gpc observed; \texttt{analysis/model\_gap\_report.md}). Two further symptoms of the same limitation: the high-cadence mode had to be encoded by conflating the pipeline requirement $i$ with the schedule ($i=6$ ``because the survey takes 6 visits''), and the cadence band $t_{\rm night}/i\le t_{\rm cad}<1$\,day had to be declared forbidden.

The extension below replaces the ``one gap'' assumption by an explicit periodic \emph{visit schedule} and rederives each of the five roles from it. Every formula reduces identically to the legacy model at $N_v=1$, and the analytic (dominant-term) mode remains a closed-form piecewise model built from the existing A1--A7 machinery.

\subsection{Schedule definition and notation}

In optical-survey mode, on the discrete cadence band ($t_{\rm cad}=n\cdot86400$\,s, $n\in\mathbb N$), each field is observed on every $n$-th night (its \emph{visit-night}) with $N_v$ visits at uniform spacing $\Delta t_v$, all inside the night window $t_{\rm night}$. Relative to the field's schedule, the visit times in one period are
\[
v_k = k\,\Delta t_v,\qquad k=0,\dots,N_v-1,
\]
repeating with period $t_{\rm cad}$. The gap \emph{preceding} visit $k$ is
\[
g_k=\begin{cases}
G\equiv t_{\rm cad}-(N_v-1)\,\Delta t_v & k=0\quad\text{(the overnight gap, back to the previous visit-night)}\\
\Delta t_v & k=1,\dots,N_v-1 ,
\end{cases}
\]
so that $\sum_k g_k=t_{\rm cad}$ exactly --- the gaps tile the whole period, which is what later makes the phase probabilities sum to one.

\paragraph{Worked timeline (ZTF public).} $n=2$ nights, $N_v=2$, $\Delta t_v=2$\,h. One period of a field's clock:
\begin{center}
\small
\begin{tabular}{lllll}
\toprule
$t$ [h] & 0 & 2 & 2--48 & 48 \\
event & visit $v_0$ & visit $v_1$ & \emph{nothing} (rest of night, day, off-night, day) & next $v_0$ \\
preceding gap & $G=46$\,h & $\Delta t_v=2$\,h & --- & $G=46$\,h \\
\bottomrule
\end{tabular}
\end{center}
A transient that turns on at a uniformly random time lands, with probability $46/48\approx96\%$, in the long $G$ gap --- so its first two visits are the \emph{intra-night pair} 2\,h apart; only with probability $2/48\approx4\%$ does it land between $v_0$ and $v_1$, in which case its second detection waits for the next visit-night. This 96/4 split is the whole mechanism behind the corrected (much larger) detection horizon.

Two integers characterise how $i$ required detections wrap around nights:
\[
\mu=\left\lfloor\frac{i-1}{N_v}\right\rfloor,\qquad
\rho=(i-1)\bmod N_v ,
\]
and the \emph{span} of $i$ consecutive visits starting from visit $k$, $S_k=v_{k+i-1}-v_k$ (periodic extension $v_{k+mN_v}=v_k+m\,t_{\rm cad}$), answers the question ``if my first detection is visit $k$, how long until my $i$-th?''.  It takes only two values:
\begin{equation}
\label{eq:spans}
S_k=\begin{cases}
S_A\equiv\mu\,t_{\rm cad}+\rho\,\Delta t_v & k\le N_v-1-\rho\quad(N_v-\rho\text{ starting visits, incl.\ }k=0)\\[2pt]
S_B\equiv(\mu{+}1)\,t_{\rm cad}-(N_v-\rho)\,\Delta t_v & k\ge N_v-\rho\quad(\rho\text{ starting visits}).
\end{cases}
\end{equation}
(For $k+\rho\le N_v-1$ the $(i{-}1)$-th subsequent visit lands $\mu$ nights later in slot $k+\rho$, giving $S_A$ independent of $k$; otherwise it wraps one extra night, giving $S_B=S_A+G-\Delta t_v$.) $S_A$ is the ``intra-anchored'' span --- for $i\le N_v$ it is simply $(i-1)\,\Delta t_v$, the intra-night run --- while $S_B$ straddles one extra overnight gap.

The identification requirement $i$ is a property of the detection \emph{pipeline} (how many epochs it needs) and is now decoupled from the schedule's $N_v$ (how many epochs the survey offers per night): ZTF high-cadence is $(N_v=6,\ i=2)$, not $i=6$.

The schedule replaces, rather than coexists with, sub-day cadences: in optical mode only $t_{\rm cad}=n\cdot86400$\,s is valid (Sec.~\ref{sec:optical_constraints}), since revisiting within the night is precisely what $(N_v,\Delta t_v)$ parameterise on the $n=1$ row. Non-optical mode (no night structure) keeps the original continuous-cadence model and ignores $(N_v,\Delta t_v)$.

\subsection{Unified exposure-time budget and $f_{\rm eff}$}
\label{sec:sched_budget}

\paragraph{Motivation for replacing $f_{\rm live}$.} The wall-clock live fraction $f_{\rm live}$ conflates two unrelated things: the astronomical night fraction $f_{\rm night}=t_{\rm night}/86400$ (geometry) and how much of the night the program's exposures actually occupy (its time share of the telescope, plus engineering losses). Define instead
\[
f_{\rm eff}\in(0,1]:\quad\text{usable fraction of the night window},\qquad f_{\rm live}=f_{\rm eff}\,f_{\rm night}.
\]
$f_{\rm eff}$ budgets telescope \emph{time} only.  Which field-nights yield usable \emph{discovery} data --- weather, moon/background quality, Galactic-plane and reference-image losses --- is a different loss and is carried by the external coverage factor $\varepsilon_{\rm cov}$ of the validation analysis (Sec.~\ref{sec:sched_comparison}); charging weather to both would double-count it.  A useful diagnostic follows: \emph{different observing modes of the same telescope should have similar $f_{\rm eff}$} (same telescope, same nights) --- a large mismatch signals a modelling error. The legacy calibration indeed carried one: ZTF public and high-cadence needed $f_{\rm live}=0.08$ vs $0.17$ ($\times2.1$ apart), an artefact of the public mode's hidden second visit per night (below). The replacement is optical-only: non-optical mode has no night window, so $f_{\rm eff}$ is undefined there and the wall-clock $f_{\rm live}$ remains that mode's own budget parameter.

\paragraph{Budget derivation.} The survey footprint is $\Omega_{\rm srv}=N_{\rm exp}\Omega_{\rm exp}$. Fields are staggered so that each night $N_{\rm exp}/n$ fields are due, each receiving $N_v$ visits. The program's usable time per night, $f_{\rm eff}\,t_{\rm night}$, must accommodate them: $(N_{\rm exp}/n)\,N_v\,(t_{\rm exp}+t_{\rm OH})=f_{\rm eff}\,t_{\rm night}$, i.e.
\begin{equation}
\label{eq:budget_unified}
t_{\rm exp}=\frac{f_{\rm eff}\,f_{\rm night}\,t_{\rm cad}}{N_{\rm exp}\,N_v}-t_{\rm OH},
\end{equation}
using $t_{\rm cad}=n\cdot86400$ and $f_{\rm night}\cdot86400=t_{\rm night}$. Per-period accounting ($n\,f_{\rm eff}\,t_{\rm night}$ of time for $N_{\rm exp}N_v$ exposures) gives the identical equation, so the staggering assumption adds no constraint beyond schedulability (Sec.~\ref{sec:sched_feasibility}).

\paragraph{Exact correspondence with the legacy branches.}
(i)~$N_v=1$: Eq.~\eqref{eq:budget_unified} is the legacy discrete-band budget $t_{\rm exp}=f_{\rm live}t_{\rm cad}/N_{\rm exp}-t_{\rm OH}$ with $f_{\rm live}=f_{\rm eff}f_{\rm night}$.
(ii)~The legacy \emph{continuous} branch (nightly revisits at spacing $\Delta t<t_{\rm night}$) is the degenerate schedule $n=1$, $N_v=t_{\rm night}/\Delta t$: Eq.~\eqref{eq:budget_unified} gives
$t_{\rm exp}=f_{\rm eff}f_{\rm night}\cdot86400/(N_{\rm exp}N_v)-t_{\rm OH}=f_{\rm eff}\,\Delta t/N_{\rm exp}-t_{\rm OH}$
--- exactly the legacy $(f_{\rm live}/f_{\rm night})\,t_{\rm cad}/N_{\rm exp}-t_{\rm OH}$. The schedule therefore \emph{subsumes} the continuous branch, which is removed (Sec.~\ref{sec:optical_constraints}); both legacy model instances (\texttt{model\_day}, \texttt{model\_night}) collapse into one model, and the legacy physicality constraint $f_{\rm night}\ge f_{\rm live}$ becomes the trivial $f_{\rm eff}\le1$ (the dynamic $t_{\rm night}$-slider floor is retired).

\paragraph{ZTF consistency check.} Anchoring both ZTF modes to their real $30$\,s exposures ($+15$\,s overhead), Eq.~\eqref{eq:budget_unified} gives
\begin{center}
\small
\begin{tabular}{lccccc}
\toprule
Mode & $N_{\rm exp}$ & $t_{\rm cad}$ & $N_v$ & exposures/night & $f_{\rm eff}$ \\
\midrule
public ($15{,}000$\,deg$^2$) & $15000/47\approx319$ & 2\,d & 2 & $\approx319$ & $319\times45\,\mathrm s/36000\,\mathrm s\approx0.40$ \\
high-cadence ($2{,}500$\,deg$^2$) & $2500/47\approx53$ & 1\,d & 6 & $\approx319$ & $\approx0.40$ \\
\bottomrule
\end{tabular}
\end{center}
Both modes land on $f_{\rm eff}\approx0.40$ --- equal, as the diagnostic demands, and consistent with ZTF's documented $\approx40\%$ public time share. The legacy $\times2.1$ split in $f_{\rm live}$ was exactly the public mode's unmodelled second visit.

\subsection{Detection probability $P_i(T)$}
\label{sec:sched_Pi}

\paragraph{Intuition.} Imagine sliding a window of fixed length $T$ (the time the burst stays above threshold) along the periodic visit timeline, with a uniformly random starting phase. $P_i(T)$ is simply the fraction of starting phases for which the window covers at least $i$ visits. Because the visit pattern repeats every $t_{\rm cad}$, this fraction is a sum of ``passing lengths'', one per gap: a window starting inside gap $k$ catches visit $k$ first, and it catches $i$ visits precisely when it starts late enough in the gap that the $i$-th visit still fits inside it. Each gap therefore contributes a piece that grows linearly with $T$ (from zero once $T$ exceeds the span $S_k$, saturating at the full gap length $g_k$), and dividing by the period turns lengths into probabilities.

Formally: a burst at $(q,D)$ is detectable over a window of length $T(q,D)$ (Sec.~\ref{sec:sched_window}). The burst's phase relative to the schedule is uniform, $u\sim\mathcal U(0,t_{\rm cad})$. The window $[u,\,u+T]$ contains at least $i$ visits iff, conditioning on the first visit at or after $u$ being visit $k$ (i.e.\ $u$ falls in the gap of length $g_k$ preceding it), the $i$-th visit $v_{k+i-1}$ still fits: $v_{k+i-1}\le u+T\iff u\ge v_{k+i-1}-T$. The passing measure inside gap $k$ is $\min(g_k,\,[T-S_k]_+)$, so
\begin{equation}
\label{eq:P_i}
P_i(T)=\frac{1}{t_{\rm cad}}\Bigl[\min\bigl(G,\,[T-S_A]_+\bigr)
+(N_v-1-\rho)\,\min\bigl(\Delta t_v,\,[T-S_A]_+\bigr)
+\rho\,\min\bigl(\Delta t_v,\,[T-S_B]_+\bigr)\Bigr],
\end{equation}
a piecewise-linear function of $T$ with at most four breakpoints. It replaces, in one object, the legacy window criterion, the $u$-uniform random-start clip, \emph{and} the $\times f_{\rm night}$ factor.

\paragraph{Limits and checks.}
\begin{itemize}
\item \textbf{$N_v=1$} ($\rho=0$, $\mu=i-1$, $G=t_{\rm cad}$): $P_i(T)=\operatorname{clip}\!\bigl(T/t_{\rm cad}-(i-1),\,0,\,1\bigr)$ --- the uniform-cadence ramp whose support endpoints $(i-1)\,t_{\rm cad}$ and $i\,t_{\rm cad}$ are precisely the two legacy window conventions (\texttt{win\_i\_minus\_one} on/off, Sec.~\ref{sec:window_settings}).
\item \textbf{Saturation}: $P_i=1$ exactly at $T_{100}=\max_c\,(S_c+g_c)$ over the \emph{nonempty} channels of Sec.~\ref{sec:sched_dominant}. Channel $A_G$ gives $S_A+G$ and channel $B$'s $S_B+\Delta t_v=S_A+G$ identically, so $T_{100}=S_A+\max(G,\Delta t_v)$ in general but $S_A+G$ whenever the $A_\Delta$ channel is empty ($\rho=N_v-1$ --- including every $N_v=1$ case, where $T_{100}=i\,t_{\rm cad}$ regardless of the inert $\Delta t_v$). \texttt{T\_certain()} implements exactly this.
\item \textbf{$\Delta t_v\to0$}: $P_i\to\operatorname{clip}\!\bigl(T/t_{\rm cad}-\mu,\,0,\,1\bigr)$ --- $N_v$ simultaneous visits act as one, and only the $\lceil i/N_v\rceil$-night accumulation survives ($\mu=\lceil i/N_v\rceil-1$).
\item \textbf{Short events, $i\le N_v$} ($\mu=0$, $S_A=(i-1)\Delta t_v$): for $T$ just above $S_A$ the first two terms ramp together, $P_i\approx(N_v-i+1)\,(T-S_A)/t_{\rm cad}$, saturating near $(N_v-i+1)\,\Delta t_v/t_{\rm cad}$. On the degenerate nightly schedule ($t_{\rm cad}=1$\,d, $N_v=t_{\rm night}/\Delta t_v$, $i\ll N_v$) this is $\approx t_{\rm night}/86400=f_{\rm night}$ --- the legacy continuous-branch night factor, now derived rather than imposed.
\item \textbf{Monotonicity}: $P_i(T)$ is non-decreasing in $T$ and non-increasing in $i$; $\sum_k g_k=t_{\rm cad}$ guarantees $P_i\to1$ as $T\to\infty$.
\end{itemize}

\paragraph{Worked example --- ZTF public} ($i=2$, $N_v=2$, $t_{\rm cad}=2$\,d; $\mu=0$, $\rho=1$, $S_A=\Delta t_v$, $S_B=2\,\mathrm d-\Delta t_v$, and the middle term of Eq.~\eqref{eq:P_i} is empty):
\[
P_2(T)=\frac{\min\bigl(G,[T-\Delta t_v]_+\bigr)+\min\bigl(\Delta t_v,[T-(2\,\mathrm d-\Delta t_v)]_+\bigr)}{2\,\mathrm d}.
\]
Reading: with probability $G/t_{\rm cad}\approx96\%$ the first two visits after the burst are an \emph{intra-night pair} spanning only $\Delta t_v$; with probability $\Delta t_v/t_{\rm cad}$ the burst lands between the pair and the two detections straddle the overnight gap. An event visible for $T\sim$\,half a period is caught with $P\approx T/t_{\rm cad}$ --- the ``factor $1/n$ for the missed nights''. The legacy encoding demanded $T_{\rm req}=4$\,d with certainty; the corrected requirement is a few hours at $96\%$ weight. Numerically, with $\Delta t_v=2$\,h ($S_A=2$\,h, $G=S_B=46$\,h):
\begin{center}
\small
\begin{tabular}{lcccccc}
\toprule
$T$ & 2\,h & 6\,h & 12\,h & 24\,h & 46\,h & 48\,h \\
$P_2(T)$ & 0 & 8\% & 21\% & 46\% & 92\% & 100\% \\
\bottomrule
\end{tabular}
\end{center}
--- a 6-hour transient, invisible to the legacy 4-day requirement, is already caught 8\% of the time; certainty needs the full period $S_A+G=48$\,h.

\subsection{Detection window $T(q,D)$}
\label{sec:sched_window}

$P_i$ is a function of a \emph{window length}, so the detection criterion must be a window condition. The legacy from-burst rectangle ($t_p\le T_{\rm req}$ \textbf{and} $t_+\ge T_{\rm req}$) is not one --- it credits the pre-peak interval where the off-axis flux is negligible (Sec.~\ref{sec:window_settings}) --- but the corrected criterion of Eq.~\eqref{eq:window_corrected} is: the above-threshold window is $[\,t_{p,\rm eff}(q),\,t_+(q,D)\,]$ with
\[
T(q,D)=t_+(q,D)-t_{p,\rm eff}(q).
\]
The crossing time $t_+$ is the \textbf{two-phase} inversion of the post-peak decline --- the same $D(t_+)$ relation as Eq.~\eqref{eq:D_window}, read in the other direction:
\begin{equation}
\label{eq:t_plus_two_phase}
t_+(q,\tilde D)=
\begin{cases}
t_{p,\rm eff}\,(\tilde D_{\max}/\tilde D)^{k}, & \tilde D\ge\tilde D_b,\\[4pt]
t_j\,(\tilde D_b/\tilde D)^{k_{\rm III}}, & \tilde D<\tilde D_b,
\end{cases}
\qquad
\tilde D_b\equiv\tilde D_{\max}\,(t_{p,\rm eff}/t_j)^{1/k},
\end{equation}
with $k=2/|\alpha_\Phi|$ phase-selected at the peak and $\tilde D_b$ the distance whose crossing lands exactly at the jet break: closer events stay visible past $t_j$, where the decline steepens to $\alpha_{\rm III}$, so extrapolating the shallow pre-break slope would overestimate their windows (by up to a factor $\sim\!2$ in rate at multi-day cadences or $i\ge3$; at the ZTF presets the correction is small --- $1.4\%$ public, $0.1\%$ high-cadence).  Phase-III viewers ($t_{p,\rm eff}\ge t_j$) are single-phase exact already (\texttt{\_t\_plus\_sched()}, matching the legacy $N_v=1$ from-peak window \texttt{D\_from\_t\_plus()} --- the two modes share one inversion).  Every visit inside the window is above the flux limit, so counting visits is counting detections. The schedule machinery therefore composes with the existing window settings per channel (Sec.~\ref{sec:sched_dominant}): \texttt{win\_from\_peak} keeps its meaning (window measured from the peak), and \texttt{win\_i\_minus\_one} selects the optimistic span convention where a hard endpoint window survives.

\subsection{Dominant-term (analytic) mode: gap-averaged mixture of rectangles}
\label{sec:sched_dominant}

Group the schedule phases by the gap the burst lands in --- intuitively: ``which kind of wait did this burst draw?''.  There are at most three \emph{channels} $c$, each a group of identical starting gaps:
\begin{center}
\small
\renewcommand{\arraystretch}{1.3}
\begin{tabular}{lcccl}
\toprule
Channel & multiplicity $m_c$ & gap $g_c$ & span $S_c$ & meaning \\
\midrule
$A_G$ & $1$ & $G$ & $S_A$ & first detection opens a visit-night \\
$A_\Delta$ & $N_v-1-\rho$ & $\Delta t_v$ & $S_A$ & first detection mid-night, no extra wrap \\
$B$ & $\rho$ & $\Delta t_v$ & $S_B$ & the $i$ detections straddle one extra night \\
\bottomrule
\end{tabular}
\end{center}
with phase weights $w_c=m_c\,g_c/t_{\rm cad}$, $\sum_c w_c=1$. Within channel $c$ the wait $u$ from the window start to the first visit is uniform on $[0,g_c]$, and \emph{a fixed wait makes the detection criterion a hard rectangle} at duration $T=S_c+u$ (the window must cover the wait plus the span). The channel rate is therefore, exactly,
\begin{equation}
\label{eq:R_mixture}
R_{\rm dom}=\sum_c w_c\,\bar R_c,
\qquad
\bar R_c=\frac1{g_c}\int_0^{g_c} R_{\rm rect}\bigl(T=S_c+u\bigr)\,du,
\end{equation}
where $R_{\rm rect}(T)$ is the unmodified A1--A7 machinery with $q_i(T)$, $D_i(T)$ and the per-channel cut baselines (Sec.~\ref{sec:sched_cuts}) substituted for their $i\,t_{\rm cad}$ forms. $R_{\rm rect}$ is piecewise power-law in $T$, so a fixed 5-point Gauss--Legendre rule evaluates the average (quadrature error $\lesssim2\%$ --- far below the endpoint spread it replaces, below). Under \texttt{win\_from\_peak} no closed rectangle exists in $T$ and the rate falls back to the exact $q$-integral, exactly mirroring the legacy fallback.

\paragraph{Reduction and bracketing.} At $N_v=1$ there is a single channel and the legacy model is used verbatim (that case never reaches this code path), with its live window toggle $T_{\rm req}=i\,t_{\rm cad}$ / $(i-1)\,t_{\rm cad}$. At $N_v\ge2$ the endpoint conventions $T=S_c+g_c$ (worst-case wait) and $T=S_c$ (zero wait) are exact bounds of the average, because
\[
\sum_c w_c\,\mathbf 1\{T\ge S_c+g_c\}\;\le\;P_i(T)\;\le\;\sum_c w_c\,\mathbf 1\{T\ge S_c\}
\]
pointwise --- the average of Eq.~\eqref{eq:R_mixture} \emph{resolves} this bracket the same way the exact mode's phase ramp does, so the \texttt{win\_i\_minus\_one} toggle is inert here.  (Earlier iterations used the endpoints directly; at the ZTF public preset they straddled the exact mode by $\times0.3$ to $\times3.7$, while the gap average lands within $\sim$10\% of it with the cuts off.)

\paragraph{Worked example --- ZTF public, continued.} With $\rho=1$, $N_v=2$ the $A_\Delta$ channel is empty. The conservative endpoints of both surviving channels coincide at $S_A+G=S_B+\Delta t_v=t_{\rm cad}=2$\,d --- already a factor-2 shorter requirement than the legacy $i\,t_{\rm cad}=4$\,d, because only \emph{one} inter-night wait is needed, not two --- while the optimistic endpoints are dominated by the intra-night pair at $S_A=\Delta t_v$. The gap average integrates $R_{\rm rect}$ across each bracket: numerically $97/{\rm yr}$ against the exact mode's $92/{\rm yr}$ at the preset (cuts off), where the endpoints gave $29$ and $349/{\rm yr}$.

\paragraph{Regime classification and medians.} The displayed regime (colouring/hover only) is classified with the single endpoint-convention duration --- $T_{100}$ of \texttt{T\_certain()} (legacy convention) or $\min_c S_c$ of \texttt{T\_first\_pass()} (\texttt{win\_i\_minus\_one}) --- and the $(S_A,G)$ cut baselines; at $N_v=1$ this is the unchanged legacy classification. The detected population of Eq.~\eqref{eq:R_mixture} is the (channel, Gauss-node) mixture of rectangle populations, so the closed-form medians of Eq.~\eqref{eq:filter_medians} no longer apply for $N_v\ge2$: \texttt{compute\_medians} takes the numerical path (the forcing already used under \texttt{win\_from\_peak}) with the \emph{same} node mixture as the rate --- guaranteeing that every cell with a drawn rate has finite medians (the hover invariant). Channels whose rectangle rate is zero simply drop out of the sum; a cell is invalid only when every channel is.

\subsection{Identification filters per channel}
\label{sec:sched_cuts}

Both identification cuts acquire the channel's local timescales --- no new formulas, only correct baselines:
\begin{itemize}
\item \textbf{Fade baseline}: the $i$ detections span $S_c$, so Eq.~\eqref{eq:tfs_discrete} is evaluated with $(i-1)\,t_{\rm cad}\to S_c$. For the intra-anchored channels of a short-$i$ survey ($S_A=(i-1)\Delta t_v$, hours) the measured decline is tiny and the cut is weak --- honestly: little fades in two hours.
\item \textbf{Rise gap}: the last non-detection before the first detection is the visit one gap earlier, so Eq.~\eqref{eq:eta_rise} is evaluated with $t_{\rm cad}\to g_c$. Channel $A_G$ --- the dominant one, weight $G/t_{\rm cad}$ --- gets the \emph{overnight} gap $G$ and hence a strong, honest cut (``absent last visit-night''); the $\Delta t_v$-gap channels get $\eta\approx1$, correctly: a source absent $\Delta t_v$ earlier is a spectacular riser and passes any rise cut trivially. This repairs the legacy pathology (noted in \texttt{analysis/model\_gap\_report.md}) in which \emph{all} events at short cadence received the toothless short gap.
\end{itemize}
The per-cut hard-boundary switches (\texttt{rise\_random\_start}, \texttt{fade\_random\_start}) keep their exact semantics, applied channel-wise.

\subsection{Exact mode: per-channel weighted volumes}
\label{sec:sched_exact}

Within channel $c$ the wait from window start to first visit is $u'\sim\mathcal U(0,g_c)$. Detection requires $u'\le T-S_c$; the fade cut $u'\le\tau_{f,c}\equiv t_{f,s}(S_c)-t_{p,\rm eff}$; the rise cut $u'\le\tau_{r,c}\equiv t_{\rm lim}(\eta(g_c))-t_{p,\rm eff}$ (Eq.~\eqref{eq:t_lim}). All three are conditions on the \emph{same} $u'$ (comonotone, exactly as in Eq.~\eqref{eq:w_joint}), so the channel's passing measure is a single clipped minimum and the joint schedule weight is
\begin{equation}
\label{eq:W_schedule}
W(q,\tilde D)=\frac{1}{t_{\rm cad}}\sum_c m_c\,\Bigl[\min\bigl(g_c,\;T(q,\tilde D)-S_c,\;\tau_{f,c},\;\tau_{r,c}\bigr)\Bigr]_+ .
\end{equation}
With the cuts off ($\tau\to+\infty$) Eq.~\eqref{eq:W_schedule} integrates $P_i$ of Eq.~\eqref{eq:P_i} exactly. Under \texttt{win\_from\_peak} the $T-S_c$ term implements the detection ramp with the two-phase crossing time of Eq.~\eqref{eq:t_plus_two_phase}; with it off, the detection window instead keeps the per-channel hard rectangle $D_i(T_{{\rm req},c})$ (matching the legacy exact-mode treatment, where the window is a hard cap and only the cuts ride the phase ramp). The rise deadline $t_{\rm lim}$ keeps the \emph{single-phase} inversion of Eq.~\eqref{eq:t_lim}, exactly like the legacy exact mode at $N_v=1$ --- an approximation both paths share by construction.

The full-integral rate replaces $V_w$ of Eq.~\eqref{eq:V_w} by the channel sum of the same closed-form pieces: each term of Eq.~\eqref{eq:W_schedule} is piecewise $\{$const $g_c$ $|$ ramp in $T$ $|$ ramp/const in $\tau\}$, with the detection ramp split at the jet-break distance $\tilde D_b$ of Eq.~\eqref{eq:t_plus_two_phase} into two segments of the same antiderivative family $G(\tilde D)$ (pre-break with $k$, post-break with $k_{\rm III}$ and $t_j$ as the anchor).  One corner survives outside closed form: with \emph{both} ramps active, below $\tilde D_b$ the detection ramp carries $k_{\rm III}$ while the rise ramp carries $k$, and mixed exponents have no closed-form crossing --- that low-volume-weight segment is integrated by a composite Gauss--Legendre rule in $\log v$ ($v=\tilde D^3$; the smooth pieces are power laws, i.e.\ exponentials in $\log v$), split at the ramps' closed-form clip corners ($\tau=m_0$ and $\tau=0$ distances) so each piece is smooth up to the single crossing kink.  This is in-style: the exact mode already integrates $q$ numerically.  Cost: a Python-level loop over $\le3$ channels inside the existing $q$-chunk loop, on discrete-band optical cells only. The numerical medians and the $dR/dq$, $dR/dD$ views inherit the channel sum through the same two entry points ($V_w$-analogue and pointwise $W$); the medians' distance-level grid is additionally capped by the largest per-channel window so its resolution follows the detected population, not the flux horizon.

\subsection{Feasibility constraints}
\label{sec:sched_feasibility}

Joining the validity region A0 on the discrete band (masks, not UI clamps --- infeasible cells render invalid):
\[
t_{\rm exp}>0,\qquad
(N_v-1)\,\Delta t_v\le t_{\rm night}\quad\text{(the visit run fits the night)},\qquad
\Delta t_v\ge t_{\rm exp}+t_{\rm OH}\quad\text{(no overlapping revisit)}.
\]

\subsection{Assumptions and approximations}
\begin{enumerate}
\item Fields are staggered uniformly across the $n$ nights of the period; the per-night and per-period budget accountings then coincide (Sec.~\ref{sec:sched_budget}).
\item Visits are equally spaced at $\Delta t_v$, all inside a top-hat night window, all at equal depth $F_{\rm lim}(t_{\rm exp})$.
\item Every visit inside the above-threshold window counts as a detection, so ``$\ge i$ visits inside the window'' is the criterion --- which is exactly a real pipeline's ``any $i$ epochs'' requirement (e.g.\ ZTF's 2-epoch alert rule), not a stricter consecutiveness demand. Band structure ($g$ vs $r$) is not modelled --- $i$ counts epochs of any band.
\item The detection window uses the two-phase crossing time of Eq.~\eqref{eq:t_plus_two_phase}; the rise deadline $t_{\rm lim}$ remains single-phase ($k=2/|\alpha_\Phi|$, phase-selected), exactly as in the legacy exact mode.
\item The dominant mode averages each channel's rectangle over its gap (Eq.~\eqref{eq:R_mixture}) --- the same wait integration the exact mode performs through the phase ramp --- but its identification cuts remain best-case hard boundaries, so it stays optimistic when the cuts bite (e.g.\ $\times1.9$ at the ZTF public preset with $s_{\rm fade}{=}0.3$, $s_{\rm rise}{=}0.5$). Exact mode is the reference.
\item In optical mode only integer-day cadences exist (Sec.~\ref{sec:optical_constraints}); intra-night sampling lives entirely in $(N_v,\Delta t_v)$, e.g.\ ZTF high-cadence is $(t_{\rm cad}{=}1\,\mathrm d,\ N_v{=}6)$, where the honest overnight rise gap applies. The old gap-band physics ($i$ detections not fitting one night) is expressible as $n{=}1$, $N_v<i$: $P_i$ accumulates detections across consecutive nights.
\end{enumerate}

\paragraph{Real-world effects absorbed rather than modelled (sizes).} The top-hat, equal-depth night hides several known effects that are deliberately left inside the two efficiency factors rather than modelled: (i)~\emph{band alternation} --- at the afterglow's spectral slope $\beta=-(p-1)/2\approx-0.6$ a $g$/$r$ pair differs by $\approx0.19$\,mag in afterglow colour, largely offset by the deeper $g$ limit; net $\lesssim0.2$\,mag per-visit depth asymmetry $\Rightarrow\lesssim25\%$ in rate ($R\propto F_{\rm lim}^{-3/2}$), partially cancelling between the bands. (ii)~\emph{Airmass/moon depth variation} across the night is a $\sim$10--20\% convexity correction on the equal-depth assumption. (iii)~\emph{Night-correlated weather}: a flat $\varepsilon_{\rm cov}$ is structurally correct for the intra-night channels (one clear night = one usable field-night; $\approx96\%$ of the phase weight at the public preset), and mildly optimistic only for cross-night detections --- the $A$-channels touch $\mu+1=\lceil i/N_v\rceil$ visit-nights and channel $B$ one more, so those detections need $\sim f_w^{\mu+1}$ / $f_w^{\mu+2}$ clear nights rather than one factor. All three are far below the current calibration frontier (the single-luminosity residual of Sec.~\ref{sec:sched_comparison}); revisit only if $i>N_v$ strategies are ever quoted.

\paragraph{Deferred refinements.} Optimiser search over $(N_v,\Delta t_v)$; closed-form mixture medians; per-band epochs.

\paragraph{Implementation reference.} \texttt{grb\_detect/schedule.py}: \texttt{NightSchedule($N_v$, dt\_v\_s, t\_night\_s)} (frozen), \texttt{spans()} (Eq.~\eqref{eq:spans}: $S_A$, $S_B$, $G$, multiplicities; a sub-period visit run, $G<0$, poisons the cells with NaN so direct API misuse reads invalid rather than $P_i>1$), \texttt{P\_detect()} (Eq.~\eqref{eq:P_i}), \texttt{joint\_weight()} (Eq.~\eqref{eq:W_schedule}), \texttt{T\_certain()}/\texttt{T\_first\_pass()} (the $P_i$ support endpoints), \texttt{feasible()} (Sec.~\ref{sec:sched_feasibility}). \texttt{DetectionRateModel(schedule=...)} with \texttt{\_t\_plus\_sched()} (Eq.~\eqref{eq:t_plus_two_phase}) and the 5-point gap-averaged mixture (Eq.~\eqref{eq:R_mixture}); \texttt{make\_rate\_model(f\_live, N\_v, dt\_v\_s, t\_night\_s, schedule\_on, ...)} --- note the engine takes the wall-clock \texttt{f\_live}; the $f_{\rm eff}\,f_{\rm night}$ conversion lives in the bridge. A single model instance replaces \texttt{model\_day}/\texttt{model\_night} (budget Eq.~\eqref{eq:budget_unified} internal, $f_{\rm night}$ factor internal). Bridge parameter keys: \texttt{f\_eff}, \texttt{N\_v}, \texttt{dt\_v\_h} (a legacy \texttt{f\_live}-only caller is accepted and used directly as the wall-clock fraction). Tests: \texttt{tests/test\_schedule\_pi.py} (Eq.~\eqref{eq:P_i} vs brute force), \texttt{tests/test\_schedule\_model.py} ($N_v{=}1$ parity, the gap-average identity and its bracket, two-phase consistency, budget unification, cut baselines), \texttt{tests/test\_schedule\_parity.py} (full-payload regression snapshots).

% -------------------------------------------------------
\section{Schedule Model vs.\ Pre-Schedule Model vs.\ Observations}
\label{sec:sched_comparison}

This section compares the night-schedule model of Sec.~\ref{sec:schedule} against both the pre-schedule (one-timescale) model and the ZTF observations, at the two real observing modes.  Numbers are from \texttt{analysis/ztf\_validation.py} (run of 2026-08-24, with the two-phase window of Eq.~\eqref{eq:t_plus_two_phase} --- a $1.4\%$ (public) / $0.1\%$ (high-cadence) shift at these presets; saved in \texttt{analysis/validation\_output.txt}); the observational targets and the distance-yardstick decision (comoving $D_C$ as the primary measure for a rate-calibrated Euclidean model) are documented in \texttt{analysis/model\_gap\_report.md}.

\subsection{Setup}

Both encodings are evaluated with the current calibrated physics ($\varepsilon_B=10^{-4}$, $D_{\rm euc}=4.55$\,Gpc, $p=2.2$), in exact rate mode with the from-peak window, the identification cuts $s_{\rm fade}=0.3$, $s_{\rm rise}=0.5$\,mag/day, and the external coverage factor $\varepsilon_{\rm cov}=0.35$ (weather, moon, Galactic plane, reference availability — a data-quality loss, distinct from the time budget $f_{\rm eff}$):

\begin{center}
\small
\renewcommand{\arraystretch}{1.25}
\begin{tabular}{lll}
\toprule
mode & pre-schedule encoding & schedule encoding \\
\midrule
A (public) & $i{=}2$, 1 visit / 2 nights, $f_{\rm eff}{=}0.19$ & $i{=}2$, \textbf{2 visits/night} @ 2\,h, every 2 nights, $f_{\rm eff}{=}0.40$ \\
B (high-cad) & $i{=}6$ (conflated), $\Delta t{=}t_{\rm night}/6$ cont. & \textbf{$i{=}2$}, 6 visits/night @ 1.5\,h, nightly, $f_{\rm eff}{=}0.40$ \\
\bottomrule
\end{tabular}
\end{center}

\subsection{Results}

\begin{center}
\small
\renewcommand{\arraystretch}{1.3}
\begin{tabular}{lcccccc}
\toprule
 & \multicolumn{3}{c}{Mode A (public)} & \multicolumn{3}{c}{Mode B (high-cadence)} \\
 & pre-sched & schedule & observed & pre-sched & schedule & observed \\
\midrule
rate $\times\,\varepsilon_{\rm cov}$ [yr$^{-1}$] & 4.7 & 14.4 & 2.0 & 5.0 & 13.3 & 2.5 \\
$q_{\rm med}$ & 1.17 & 1.02 & $\lesssim 1.5$ & 1.13 & 1.05 & $\lesssim 1.5$ \\
on-axis-consistent ($q<1.5$) & 78\% & 88\% & $\ge 60\%$ & 79\% & 90\% & $\ge 60\%$ \\
$D_{\rm med}$ [Gpc, $D_C$] & 1.00 & 1.91 & 3.67 & 1.82 & 2.92 & 3.88 \\
$D_{90}$ [Gpc] & 1.41 & 3.92 & --- & 2.67 & 4.25 & --- \\
\bottomrule
\end{tabular}
\end{center}

\subsection{Findings}

\paragraph{(i) The distance gap closes by half, exactly through the hypothesized mechanism.}  Crediting the public mode's intra-night pair — a required span of $S_A=\Delta t_v\approx 2$\,h at $G/t_{\rm cad}\approx 96\%$ phase weight, instead of the pre-schedule $T_{\rm req}=i\,t_{\rm cad}\approx 4$\,d — moves the cadence horizon out and with it $D_{\rm med}$ ($\times 1.9$) and $D_{90}$ ($\times 2.8$).  The model's spurious low-vs-high-cadence median split narrows from $\times 2.5$ to $\times 1.54$ (observed: $\times 1.06$), and the angular statistics ($q_{\rm med}\approx 1.0$, 88--89\% on-axis-consistent) now sit inside the observed targets.

\paragraph{(ii) The rates overshoot ($\times 5$--$7$) --- the same structural coupling seen from the other side.}  In every regime $R\propto H^3$ while $D_{\rm med}\approx 0.79\,H$ ($H$ the binding horizon), so a model with a single burst luminosity cannot move the median outward without the rate following.  The pre-schedule encoding matched the observed rates \emph{because} its horizon was artificially short — an error cancellation, not an agreement.  Removing the encoding error exposes the underlying deficiency: no calibration of the single-luminosity model can now match rates and medians simultaneously.  This is precisely the deferred luminosity-function prescription of \texttt{analysis/model\_gap\_report.md}: a lognormal $F_{\rm dec}$ spread ($\sigma\approx 1.2$--$1.5$\,mag, literature-anchored) with a joint ($\sigma$, median-brightness) refit of the rates.  Per the standing assessment-only decision, it is \emph{not} applied here.

\paragraph{(iii) The window-convention bracket collapses.}  The pre-schedule model bracketed the detection criterion between $T_{\rm req}=i\,t_{\rm cad}$ and $(i{-}1)\,t_{\rm cad}$ (a factor ~[3, 13] in Mode A's rate across the toggles).  In exact mode the schedule integrates the phase ramp $P_i(T)$ exactly, so the two conventions now coincide by construction and the associated modeling uncertainty is resolved, not bracketed.

\paragraph{(iv) Sensitivities at the final configuration (Mode A).}  The schedule knobs are weak once the ramp is exact: $\Delta t_v\in[0.5, 5]$\,h changes the rate by $<5\%$; $N_v{=}2\to 1$ at the same budget is $\times 0.68$.  The restored overnight rise gap gives the identification cuts real teeth ($s_{\rm rise}=1.0$: $\times 0.37$ --- the pre-schedule sub-day encoding had $\eta\to 1$); $i=3$ is $\times0.24$ (the two-phase window bites once the required span crosses $t_j$); $\varepsilon_B\pm 0.5$\,dex spans $\times 0.30$--$2.8$.  The Ho et al.\ \S4.1 strict intranight benchmark is unchanged ($13.4$/yr vs.\ their MC $0.52$/yr): its overshoot was never a cadence-encoding problem and carries the same luminosity-function signature.

\paragraph{Bottom line.}  The schedule model is \emph{better} on every observable it was built to fix — distances, viewing angles, and the cross-mode consistency ($f_{\rm eff}\approx 0.40$ for both modes, vs.\ the old $\times 2.1$ split in $f_{\rm live}$) — and \emph{worse} on the raw rate normalization, in a way that is now cleanly attributable: the residual pattern (rates $\times 5$--$7$ high at horizons still $\times 1.3$--$1.9$ short) is the quantified signature of the single-luminosity population, with an identified, deferred fix.  No parameter recalibration is applied without approval.

\end{document}
